\paragraph{自然扩张的周期守恒、prime-slice 精确极小性与 Smith 核增长的 Euler 定律}
沿用定理 \ref{thm:killo-layered-prime-slices-finite-depth}、推论
\ref{cor:killo-layered-necessity-minimality}、定理
\ref{thm:killo-natural-extension-branch-register}、定理
\ref{thm:killo-smith-loss-spectrum} 与推论
\ref{cor:killo-kernel-size-general-modulus} 的记号。下列结果把
“自然扩张--有限历史 prime-slice--Smith 信息亏损谱--全模数核增长”
四段对象压缩为同一条可直接调用的结论链。

\begin{theorem}[自然扩张寄存器实现保持全部周期点计数与 Artin--Mazur \texorpdfstring{$\zeta$}{zeta} 函数]\label{thm:xi-natural-extension-periodic-orbit-zeta-preservation}
设 \(T:X\twoheadrightarrow X\) 为有限纤维满射，\(X^{\leftarrow}\) 为定理
\ref{thm:killo-natural-extension-branch-register} 的自然扩张，\(\sigma\) 为其左移同构，
\[
\mathcal C:X^{\leftarrow}\xrightarrow{\sim}\widehat X,
\qquad
\widehat T:=\mathcal C\circ \sigma\circ \mathcal C^{-1}.
\]
则对每个整数 \(n\ge 1\)，映射
\[
\pi_0:\operatorname{Fix}(\sigma^n)\longrightarrow \operatorname{Fix}(T^n),
\qquad
(x_0,x_{-1},x_{-2},\dots)\longmapsto x_0
\]
为自然双射。因而也有自然双射
\[
\operatorname{Fix}(\widehat T^n)\cong \operatorname{Fix}(T^n).
\]
特别地，若这些不动点集合逐个有限，则 Artin--Mazur \(\zeta\)-函数满足严格恒等式
\[
\zeta_{\widehat T}(z)
:=
\exp\!\left(\sum_{n\ge 1}\frac{\#\operatorname{Fix}(\widehat T^n)}{n}z^n\right)
=
\exp\!\left(\sum_{n\ge 1}\frac{\#\operatorname{Fix}(T^n)}{n}z^n\right)
=:
\zeta_T(z).
\]
进一步，若两侧拓扑熵均由周期点增长率定义，则
\[
h_{\mathrm{top}}(\widehat T)=h_{\mathrm{top}}(T).
\]
\leanverified{paper\_xi\_natural\_extension\_periodic\_orbit\_zeta\_preservation}
\end{theorem}
\begin{proof}
对任意
\[
\mathbf x=(x_0,x_{-1},x_{-2},\dots)\in X^{\leftarrow},
\]
由自然扩张的定义可知
\[
\sigma^n(\mathbf x)
=
\bigl(T^n(x_0),T^{n-1}(x_0),\dots,T(x_0),x_0,x_{-1},x_{-2},\dots\bigr).
\]
故若 \(\sigma^n(\mathbf x)=\mathbf x\)，则必有
\[
T^n(x_0)=x_0,
\qquad
x_{-r}=T^{n-r}(x_0)\quad(1\le r\le n),
\]
并且由继续比较坐标得到
\[
x_{-(qn+r)}=T^{n-r}(x_0)
\qquad(q\ge 0,\ 1\le r\le n).
\]
因此 \(\mathbf x\) 由 \(x_0\in\operatorname{Fix}(T^n)\) 唯一决定。

反之，任取 \(x_0\in\operatorname{Fix}(T^n)\)，按上式定义
\[
x_{-(qn+r)}:=T^{n-r}(x_0)
\qquad(q\ge 0,\ 1\le r\le n),
\]
则由 \(T^n(x_0)=x_0\) 可直接验证
\[
T(x_{-j})=x_{-j+1}\qquad(j\ge 1),
\]
故所得 \(\mathbf x\in X^{\leftarrow}\)；并且再由构造立即有
\[
\sigma^n(\mathbf x)=\mathbf x.
\]
从而 \(\pi_0\) 给出 \(\operatorname{Fix}(\sigma^n)\) 与 \(\operatorname{Fix}(T^n)\) 的双射。

由定理 \ref{thm:killo-natural-extension-branch-register}，
\[
(\widehat X,\widehat T)\cong (X^{\leftarrow},\sigma),
\]
因此
\[
\operatorname{Fix}(\widehat T^n)=\mathcal C\bigl(\operatorname{Fix}(\sigma^n)\bigr)
\cong
\operatorname{Fix}(T^n).
\]
若各不动点集合有限，则两侧 Artin--Mazur \(\zeta\)-函数的指数级数逐项相同，故 \(\zeta_{\widehat T}=\zeta_T\)。若拓扑熵由
\[
h_{\mathrm{top}}(S)
=
\limsup_{n\to\infty}\frac1n\log \#\operatorname{Fix}(S^n)
\]
给出，则由于对应周期点计数对一切 \(n\) 完全一致，立得
\(
h_{\mathrm{top}}(\widehat T)=h_{\mathrm{top}}(T)
\)。
证毕。
\end{proof}

\begin{theorem}[有限层 prime-slice 编码的非平凡片数精确极小性]\label{thm:xi-prime-slice-nontrivial-layer-exact-minimality}
\leanverified{paper\_xi\_prime\_slice\_nontrivial\_layer\_exact\_minimality}
设
\[
X_0\xrightarrow{\pi_0}X_1\xrightarrow{\pi_1}\cdots\xrightarrow{\pi_{k-1}}X_k
\]
为有限纤维满射塔。定义非平凡分支层集合
\[
J:=
\left\{
\,0\le j\le k-1:\ \exists\,y\in X_{j+1}\ \text{使}\ \#\pi_j^{-1}(y)\ge 2
\right\}.
\]
考虑任意按层分解的 prime-slice 型映射
\[
\Psi(x_0)=\bigl(x_k,H(x_0)\bigr),
\qquad
H(x_0)=\prod_{j=0}^{k-1}h_j(x_0),
\]
其中每个 \(h_j(x_0)\) 的全部素因子均落在预先指定且两两不交的素数片
\(\mathcal P_j\)，并假设 \(\Psi\) 为单射。定义其非平凡片数
\[
N_{\mathrm{slice}}(\Psi):=\#\{\,j:\ h_j\not\equiv 1\,\}.
\]
则必有严格下界
\[
N_{\mathrm{slice}}(\Psi)\ge |J|.
\]
反之，存在一条这类单射 \(\Psi\) 使得
\[
N_{\mathrm{slice}}(\Psi)=|J|.
\]
故 prime-slice 编码所需的最小有效片数恰等于非平凡分支层数 \(|J|\)。
\end{theorem}
\begin{proof}
先证下界。若 \(j\in J\)，则存在 \(y_j\in X_{j+1}\) 使
\[
\#\pi_j^{-1}(y_j)\ge 2.
\]
取不同点
\[
u_j\neq v_j\in \pi_j^{-1}(y_j).
\]
再固定其余层的选择，并沿塔向下选取兼容前史、向上取公共像，得到两点
\[
x_0',x_0''\in X_0
\]
使它们在第 \(j\) 层的分支选择不同，而在其余各层的对应片语义下保持相同的选取，并且终端宏态同为某个 \(x_k\in X_k\)。
若 \(h_j\equiv 1\)，则第 \(j\) 层对乘积账本不贡献任何区分信息；由于其余层的选择已固定，分层分离便给出
\[
H(x_0')=H(x_0'').
\]
于是
\[
\Psi(x_0')=\Psi(x_0''),
\]
与单射性矛盾。故对每个 \(j\in J\) 都必须有
\[
h_j\not\equiv 1.
\]
因此每个 \(j\in J\) 都必须贡献一个非平凡片，故
\[
N_{\mathrm{slice}}(\Psi)\ge |J|.
\]

再证上界。把 \(J\) 写成严格递增序列
\[
J=\{j_1<j_2<\cdots<j_r\},
\qquad r=|J|.
\]
若 \(j\notin J\)，则 \(\pi_j\) 的每条纤维都是单点；由于 \(\pi_j\) 还是满射，故 \(\pi_j\) 实为双射，从而
\[
x_j=\pi_j^{-1}(x_{j+1})
\]
可由上一层唯一恢复，无需任何额外 prime-slice 信息。

对每个 \(j_\nu\in J\)，沿用定理 \ref{thm:killo-layered-prime-slices-finite-depth} 的构造，取单射
\[
\alpha_{j_\nu}:X_{j_\nu+1}\times\NN\hookrightarrow \mathcal P_{j_\nu}
\]
及纤维枚举
\[
\iota_{j_\nu,y}:\pi_{j_\nu}^{-1}(y)\hookrightarrow \NN.
\]
对 \(x_0\in X_0\) 记其各层像为 \(x_j\in X_j\)，并定义
\[
h_{j_\nu}(x_0):=
\alpha_{j_\nu}\!\bigl(x_{j_\nu+1},\iota_{j_\nu,x_{j_\nu+1}}(x_{j_\nu})\bigr),
\qquad
h_j(x_0):=1\quad(j\notin J).
\]
置
\[
H(x_0):=\prod_{j=0}^{k-1}h_j(x_0)
=
\prod_{\nu=1}^{r}h_{j_\nu}(x_0),
\qquad
\Psi(x_0):=(x_k,H(x_0)).
\]
显然 \(N_{\mathrm{slice}}(\Psi)=r=|J|\)。

现说明 \(\Psi\) 单射。给定 \((x_k,H)\) 属于其像。由素数片两两不交与唯一分解，可从 \(H\) 中逐片抽出全部非平凡因子
\[
h_{j_1},\dots,h_{j_r}.
\]
从顶层 \(x_k\) 开始向下恢复。若当前层 \(j\notin J\)，则
\[
x_j=\pi_j^{-1}(x_{j+1})
\]
唯一确定；若 \(j=j_\nu\in J\)，则由 \(x_{j_\nu+1}\) 与素数因子 \(h_{j_\nu}\) 经
\(\alpha_{j_\nu}\) 的像上反演恢复唯一索引，再经
\(\iota_{j_\nu,x_{j_\nu+1}}\) 的像上反演恢复唯一 \(x_{j_\nu}\in\pi_{j_\nu}^{-1}(x_{j_\nu+1})\)。
如此逐层下推，唯一恢复 \(x_0\)。故 \(\Psi\) 为单射。

综上，最小非平凡片数恰等于 \(|J|\)。证毕。
\end{proof}

\begin{theorem}[Smith 信息亏损谱的离散曲率原子精确等于 \texorpdfstring{$p$}{p}-初等因子重数]\label{thm:xi-smith-loss-discrete-curvature-atoms}
\leanverified{paper\_xi\_smith\_loss\_discrete\_curvature\_atoms}
设 \(A\in M_{m\times n}(\ZZ)\) 且 \(\operatorname{rank}_{\QQ}(A)=m\)。对固定素数 \(p\)，沿用定理
\ref{thm:killo-smith-loss-spectrum} 的记号，令
\[
e_i(p):=v_p(d_i),
\qquad
f_p(k):=\sum_{i=1}^{m}\min(k,e_i(p))
\qquad(k\ge 1),
\]
并约定 \(f_p(0):=0\)。则：
\begin{enumerate}
\item \(f_p:\ZZ_{\ge 0}\to\ZZ_{\ge 0}\) 为离散凹函数，即对每个 \(k\ge 1\) 都有
\[
f_p(k+1)-2f_p(k)+f_p(k-1)\le 0.
\]
\item 对每个 \(k\ge 1\)，其二阶差分的负值恰等于赋值 \(k\) 的初等因子重数：
\[
-\bigl(f_p(k+1)-2f_p(k)+f_p(k-1)\bigr)
=
\#\{i:\ e_i(p)=k\}.
\]
\end{enumerate}
因而 \(f_p\) 的整个图像可以直接视为 \(\{e_i(p)\}\) 的离散 Newton 图，而无需额外反演算法。
\end{theorem}
\begin{proof}
由定理 \ref{thm:killo-smith-loss-spectrum}，
\[
\Delta_p(k):=f_p(k)-f_p(k-1)=\#\{i:\ e_i(p)\ge k\}
\qquad(k\ge 1).
\]
由于右端随 \(k\) 单调不增，故
\[
\Delta_p(k+1)\le \Delta_p(k),
\]
亦即
\[
f_p(k+1)-f_p(k)\le f_p(k)-f_p(k-1).
\]
这等价于
\[
f_p(k+1)-2f_p(k)+f_p(k-1)\le 0,
\]
从而 \(f_p\) 离散凹。

再对 \(\Delta_p\) 作一次差分，得到
\[
\Delta_p(k)-\Delta_p(k+1)
=
\#\{i:\ e_i(p)\ge k\}-\#\{i:\ e_i(p)\ge k+1\}
=
\#\{i:\ e_i(p)=k\}.
\]
另一方面，
\[
\Delta_p(k)-\Delta_p(k+1)
=
\bigl(f_p(k)-f_p(k-1)\bigr)-\bigl(f_p(k+1)-f_p(k)\bigr)
\]
\[
=
-\bigl(f_p(k+1)-2f_p(k)+f_p(k-1)\bigr).
\]
比较两式即得
\[
-\bigl(f_p(k+1)-2f_p(k)+f_p(k-1)\bigr)
=
\#\{i:\ e_i(p)=k\}.
\]
证毕。
\end{proof}

\begin{theorem}[模数核大小 Dirichlet 生成函数的 Euler 因子分解与单极点定律]\label{thm:xi-smith-kernel-dirichlet-euler-single-pole}
\leanverified{paper\_xi\_smith\_kernel\_dirichlet\_euler\_single\_pole}
沿用推论 \ref{cor:killo-kernel-size-general-modulus} 的设定。定义
\[
\mathcal Z_A(s):=\sum_{b\ge 1}\frac{\#\ker(A_b)}{b^s},
\qquad
\Delta_A:=\prod_{i=1}^{m}d_i.
\]
对每个素数 \(p\)，记
\[
e_i(p):=v_p(d_i),
\qquad
Z_{A,p}(s):=
\sum_{k\ge 0}p^{\,k(n-m-s)+\sum_{i=1}^{m}\min(k,e_i(p))}.
\]
则成立：
\begin{enumerate}
\item \(\mathcal Z_A(s)\) 为积性 Dirichlet 级数，并具有 Euler 因子分解
\[
\mathcal Z_A(s)=\prod_{p}Z_{A,p}(s)
=
\zeta\!\bigl(s-(n-m)\bigr)
\prod_{p\mid \Delta_A}\bigl(1-p^{\,n-m-s}\bigr)Z_{A,p}(s).
\]
\item 对每个固定素数 \(p\)，局部因子 \(Z_{A,p}(s)\) 是 \(p^{-s}\) 的有理函数。若按非降序写
\[
0=:e_0(p)\le e_1(p)\le\cdots\le e_m(p),
\qquad
E_j(p):=\sum_{i=1}^{j}e_i(p),
\]
则有显式闭式
\[
Z_{A,p}(s)
=
\sum_{j=0}^{m-1}
p^{E_j(p)}
\frac{
p^{e_j(p)(n-j-s)}-p^{e_{j+1}(p)(n-j-s)}
}{
1-p^{\,n-j-s}
}
\]
\[
\qquad\qquad
+\;
p^{E_m(p)}
\frac{
p^{e_m(p)(n-m-s)}
}{
1-p^{\,n-m-s}
}.
\]
\item \(\mathcal Z_A(s)\) 在半平面 \(\Re(s)>n-m\) 上亚纯，并且在
\[
s_0:=n-m+1
\]
处具有唯一简单极点，其留数为
\[
\Res_{s=s_0}\mathcal Z_A(s)
=
\prod_{p\mid \Delta_A}\Bigl(1-\frac1p\Bigr)Z_{A,p}(s_0).
\]
\item 因而由 Wiener--Ikehara 定理，
\[
\sum_{b\le X}\#\ker(A_b)
\sim
\frac{\mathfrak R_A}{n-m+1}\,X^{n-m+1},
\qquad
\mathfrak R_A:=\Res_{s=s_0}\mathcal Z_A(s).
\]
\end{enumerate}
\end{theorem}
\begin{proof}
由推论 \ref{cor:killo-kernel-size-general-modulus}，
\[
\#\ker(A_b)=b^{\,n-m}\prod_{i=1}^{m}\gcd(d_i,b).
\]
右端关于 \(b\) 为积性函数，因此 Dirichlet 级数 \(\mathcal Z_A(s)\) 具有 Euler 乘积分解
\[
\mathcal Z_A(s)=\prod_p Z_{A,p}(s),
\]
其中
\[
Z_{A,p}(s)
=
\sum_{k\ge 0}\frac{\#\ker(A_{p^k})}{p^{ks}}
=
\sum_{k\ge 0}p^{\,k(n-m-s)+\sum_{i=1}^{m}\min(k,e_i(p))}.
\]
这给出第 \(1\) 条的第一式。

若 \(p\nmid \Delta_A\)，则所有 \(e_i(p)=0\)，故
\[
Z_{A,p}(s)=\sum_{k\ge 0}p^{k(n-m-s)}=\frac{1}{1-p^{\,n-m-s}}.
\]
把这些“好素数”因子提出，即得
\[
\mathcal Z_A(s)
=
\zeta\!\bigl(s-(n-m)\bigr)
\prod_{p\mid \Delta_A}\bigl(1-p^{\,n-m-s}\bigr)Z_{A,p}(s),
\]
从而得到第 \(1\) 条的第二式。

现证第 \(2\) 条。对固定素数 \(p\)，把 \(k\)-求和按区间
\[
e_j(p)\le k<e_{j+1}(p)\qquad(0\le j\le m-1)
\]
以及尾区间
\[
k\ge e_m(p)
\]
分块。在第 \(j\) 个有限区间上，有
\[
\sum_{i=1}^{m}\min(k,e_i(p))
=
E_j(p)+(m-j)k,
\]
故对应求和块为
\[
\sum_{k=e_j(p)}^{e_{j+1}(p)-1}p^{\,E_j(p)}p^{\,k(n-j-s)}
=
p^{E_j(p)}
\frac{
p^{e_j(p)(n-j-s)}-p^{e_{j+1}(p)(n-j-s)}
}{
1-p^{\,n-j-s}
}.
\]
在尾区间上则有
\[
\sum_{i=1}^{m}\min(k,e_i(p))=E_m(p),
\]
因而
\[
\sum_{k\ge e_m(p)}p^{\,k(n-m-s)+E_m(p)}
=
p^{E_m(p)}
\frac{
p^{e_m(p)(n-m-s)}
}{
1-p^{\,n-m-s}
}.
\]
把这些分块相加即得所示闭式。由于它是有限和与一个几何尾项的组合，故 \(Z_{A,p}(s)\) 的确是 \(p^{-s}\) 的有理函数。

再证第 \(3\) 条。把
\[
F_A(s):=\prod_{p\mid \Delta_A}\bigl(1-p^{\,n-m-s}\bigr)Z_{A,p}(s)
\]
记作坏素数修正因子。由上式可见，对每个 \(p\mid\Delta_A\)，
\[
\bigl(1-p^{\,n-m-s}\bigr)Z_{A,p}(s)
\]
在 \(s\) 上是有限个指数函数 \(p^{-rs}\) 的线性组合，故在整个复平面上全纯。于是 \(F_A(s)\) 为有限乘积，亦为全纯函数。故
\[
\mathcal Z_A(s)=\zeta\!\bigl(s-(n-m)\bigr)F_A(s)
\]
在 \(\Re(s)>n-m\) 上的唯一极点只能来自 \(\zeta(s-(n-m))\) 在 \(s_0=n-m+1\) 处的简单极点。

在 \(s=s_0\) 处，
\[
\Res_{s=s_0}\zeta\!\bigl(s-(n-m)\bigr)=1,
\]
故
\[
\Res_{s=s_0}\mathcal Z_A(s)=F_A(s_0)
=
\prod_{p\mid \Delta_A}\Bigl(1-\frac1p\Bigr)Z_{A,p}(s_0).
\]
第 \(3\) 条成立。

最后，由 \(\#\ker(A_b)\ge 0\) 对一切 \(b\) 成立，且 \(\mathcal Z_A(s)\) 在 \(\Re(s)\ge s_0\) 上除 \(s_0\) 之外无其他极点，Wiener--Ikehara 定理适用，从而得到
\[
\sum_{b\le X}\#\ker(A_b)
\sim
\frac{\mathfrak R_A}{s_0}\,X^{s_0}
=
\frac{\mathfrak R_A}{n-m+1}\,X^{n-m+1},
\]
其中 \(\mathfrak R_A=\Res_{s=s_0}\mathcal Z_A(s)\)。证毕。
\end{proof}

\begin{theorem}[模核 Dirichlet 级数的 Euler--Stokes 尾分解与唯一极点留数]\label{thm:xi-smith-kernel-dirichlet-euler-stokes-tail-residue}
沿用定理 \ref{thm:xi-smith-kernel-dirichlet-euler-single-pole} 的记号。对每个素数 \(p\)，记
\[
\mathcal Z_{A,p}(s):=
\sum_{k\ge 0}p^{-k(s-(n-m))+\sum_{i=1}^{m}\min(k,e_i(p))},
\qquad
E_p:=\max_i e_i(p).
\]
则
\[
\boxed{\;
\mathcal Z_A(s)
=
\zeta\!\bigl(s-(n-m)\bigr)
\prod_{p\mid \Delta_A}
\Bigl(1-p^{-(s-(n-m))}\Bigr)\,
\mathcal Z_{A,p}(s).\;}
\]
并且每个坏素数局部因子都具有显式“有限前缀 + 几何尾部”分解
\[
\boxed{\;
\mathcal Z_{A,p}(s)
=
\sum_{k=0}^{E_p-1}
p^{-k(s-(n-m))+\sum_{i=1}^{m}\min(k,e_i(p))}
\;+\;
\frac{
p^{-E_p(s-(n-m))+\sum_{i=1}^{m}e_i(p)}
}{
1-p^{-(s-(n-m))}}
\;}
\]
从而 \(\mathcal Z_A(s)\) 在半平面 \(\Re(s)>n-m\) 上亚纯，并且在
\[
s_0:=n-m+1
\]
处具有唯一简单极点，其留数可写为
\[
\boxed{\;
\Res_{s=s_0}\mathcal Z_A(s)
=
\prod_{p\mid \Delta_A}
\Bigl(1-\frac1p\Bigr)
\sum_{k\ge 0}p^{-k+\sum_{i=1}^{m}\min(k,e_i(p))}.
\;}
\]
\end{theorem}
\begin{proof}
定理 \ref{thm:xi-smith-kernel-dirichlet-euler-single-pole} 已给出
\[
\mathcal Z_A(s)
=
\zeta\!\bigl(s-(n-m)\bigr)
\prod_{p\mid \Delta_A}
\Bigl(1-p^{-(s-(n-m))}\Bigr)\,
Z_{A,p}(s),
\]
其中
\[
Z_{A,p}(s)
=
\sum_{k\ge 0}p^{-k(s-(n-m))+\sum_{i=1}^{m}\min(k,e_i(p))}.
\]
故只需把 \(Z_{A,p}(s)\) 改记为 \(\mathcal Z_{A,p}(s)\)。

若 \(k\ge E_p\)，则对一切 \(i\) 都有 \(\min(k,e_i(p))=e_i(p)\)，于是
\[
\sum_{i=1}^{m}\min(k,e_i(p))=\sum_{i=1}^{m}e_i(p)
\]
稳定为常数。因此局部级数尾部成为公比
\(
p^{-(s-(n-m))}
\)
的几何级数：
\[
\sum_{k\ge E_p}
p^{-k(s-(n-m))+\sum_{i=1}^{m}e_i(p)}
=
\frac{
p^{-E_p(s-(n-m))+\sum_{i=1}^{m}e_i(p)}
}{
1-p^{-(s-(n-m))}}.
\]
加上前 \(E_p\) 项即得所示 Euler--Stokes 尾分解。

定理 \ref{thm:xi-smith-kernel-dirichlet-euler-single-pole} 同时已证明 \(\mathcal Z_A(s)\) 在 \(\Re(s)>n-m\) 上只有 \(s_0=n-m+1\) 这一唯一简单极点，并满足
\[
\Res_{s=s_0}\mathcal Z_A(s)
=
\prod_{p\mid \Delta_A}\Bigl(1-\frac1p\Bigr)\mathcal Z_{A,p}(s_0).
\]
把
\[
s_0-(n-m)=1
\]
代入 \(\mathcal Z_{A,p}(s)\) 即得
\[
\mathcal Z_{A,p}(s_0)
=
\sum_{k\ge 0}p^{-k+\sum_{i=1}^{m}\min(k,e_i(p))},
\]
从而得到所示留数公式。证毕。
\end{proof}

由此，附录中的四段对象在本链上闭合为
\[
\text{有限纤维自映射}
\Longrightarrow
\text{自然扩张寄存器实现}
\Longrightarrow
\text{周期点与 Artin--Mazur }\zeta\text{ 的严格守恒},
\]
\[
\text{Smith 核大小谱 }K_p(k)
\Longrightarrow
\text{离散曲率原子 }\#\{i:e_i(p)=k\}
\Longrightarrow
\text{Dirichlet 生成函数的 Euler 因子化}
\Longrightarrow
\text{全模数核增长的单极点定律}.
\]

\paragraph{bin-fold 稳定统计的一位塌缩与有限局部化的 \(p\)-进核刚性}
下列结果把定理 \ref{thm:fold-bin-two-state-asymptotic}、命题 \ref{prop:fold-bin-lastbit-sufficient-tv}、
定理 \ref{thm:fold-bin-renyi-divergence-limit}、推论 \ref{cor:xi-localized-integers-order-dual-quotient}、
定理 \ref{thm:cdim-star-z1s-dual-extension} 与定理
\ref{thm:conclusion-cdim-finite-rank-extension} 组织为同一条链：
可见稳定统计在指数精度上塌缩为末位两层，而可逆算术的分类信息则完全驻留在全不连通
\(p\)-进核中。

\begin{theorem}[稳定型分布的末位充分性与两层模型塌缩]\label{thm:xi-fold-lastbit-two-layer-collapse}
设
\[
X_m:=\{w\in\{0,1\}^m:\ w_iw_{i+1}=0\ \ (1\le i<m)\},
\qquad
\mu_m(w):=\frac{d_m^{\mathrm{bin}}(w)}{2^m}.
\]
定义参考分布
\[
\nu_m(w):=\varphi^{-m-w_m},
\qquad
w=(w_1,\dots,w_m)\in X_m.
\]
则 \(\nu_m\) 是 \(X_m\) 上的概率分布，并且对一切充分大的 \(m\) 与一切 \(w\in X_m\) 都有一致估计
\[
\mu_m(w)=\nu_m(w)\Bigl(1+O\!\Bigl(\Bigl(\frac{\varphi}{2}\Bigr)^m\Bigr)\Bigr).
\]
从而
\[
D_{\mathrm{TV}}(\mu_m,\nu_m)=O\!\Bigl(\Bigl(\frac{\varphi}{2}\Bigr)^m\Bigr).
\]
进一步，对每个 \(b\in\{0,1\}\) 记
\[
X_m^{(b)}:=\{w\in X_m:\ w_m=b\},
\]
并令 \(u_m^{(b)}\) 为 \(X_m^{(b)}\) 上的均匀分布。则有
\[
D_{\mathrm{TV}}\bigl(\mu_m(\,\cdot\,|X_m^{(b)}),u_m^{(b)}\bigr)
=O\!\Bigl(\Bigl(\frac{\varphi}{2}\Bigr)^m\Bigr).
\]
换言之，在稳定型层面，除末位 \(w_m\) 外的全部组合自由度都以指数速度趋于层内均匀化。
\end{theorem}
\begin{proof}
由定理 \ref{thm:fold-bin-two-state-asymptotic}，
\[
d_m^{\mathrm{bin}}(w)=\varphi^{-w_m}\Bigl(\frac{2}{\varphi}\Bigr)^m+O(1)
\]
对 \(w\in X_m\) 一致成立。两边除以 \(2^m\) 得
\[
\mu_m(w)=\varphi^{-m-w_m}+O(2^{-m}).
\]
由于
\[
\varphi^{-m-1}\le \nu_m(w)\le \varphi^{-m},
\]
故可把余项重写为
\[
\mu_m(w)=\nu_m(w)\Bigl(1+O\!\Bigl(\Bigl(\frac{\varphi}{2}\Bigr)^m\Bigr)\Bigr),
\]
且误差在 \(w\in X_m\) 上一致。

现证 \(\nu_m\) 归一。由
\[
|X_m^{(0)}|=F_{m+1},
\qquad
|X_m^{(1)}|=F_m
\]
与 Binet 公式可得恒等式
\[
F_{m+1}+\varphi^{-1}F_m=\varphi^m.
\]
因此
\[
\sum_{w\in X_m}\nu_m(w)
=\varphi^{-m}|X_m^{(0)}|+\varphi^{-m-1}|X_m^{(1)}|
=\varphi^{-m}\bigl(F_{m+1}+\varphi^{-1}F_m\bigr)=1.
\]

再由逐点相对误差估计求和，
\[
\sum_{w\in X_m}|\mu_m(w)-\nu_m(w)|
\le
C\Bigl(\frac{\varphi}{2}\Bigr)^m\sum_{w\in X_m}\nu_m(w)
=
C\Bigl(\frac{\varphi}{2}\Bigr)^m,
\]
故
\[
D_{\mathrm{TV}}(\mu_m,\nu_m)
=O\!\Bigl(\Bigl(\frac{\varphi}{2}\Bigr)^m\Bigr).
\]

最后固定 \(b\in\{0,1\}\)。由于 \(\nu_m\) 在 \(X_m^{(b)}\) 上为常数，故
\[
\nu_m(\,\cdot\,|X_m^{(b)})=u_m^{(b)}.
\]
将前述相对误差在 \(X_m^{(b)}\) 上求和得
\[
\mu_m(X_m^{(b)})=\nu_m(X_m^{(b)})\Bigl(1+O\!\Bigl(\Bigl(\frac{\varphi}{2}\Bigr)^m\Bigr)\Bigr).
\]
于是对任意 \(w\in X_m^{(b)}\)，
\[
\mu_m(w|X_m^{(b)})
=\frac{\mu_m(w)}{\mu_m(X_m^{(b)})}
=\frac{\nu_m(w)}{\nu_m(X_m^{(b)})}
\Bigl(1+O\!\Bigl(\Bigl(\frac{\varphi}{2}\Bigr)^m\Bigr)\Bigr)
=u_m^{(b)}(w)\Bigl(1+O\!\Bigl(\Bigl(\frac{\varphi}{2}\Bigr)^m\Bigr)\Bigr),
\]
且误差在 \(w\in X_m^{(b)}\) 上一致。再对 \(X_m^{(b)}\) 求和即得条件化后的总变差界。
\end{proof}

\begin{theorem}[Rényi 熵常数谱的闭式]\label{thm:xi-fold-renyi-constant-spectrum-closed-form}
对每个固定的 \(q>0\)、\(q\neq 1\)，定义
\[
H_q(\mu_m):=\frac{1}{1-q}\log\sum_{w\in X_m}\mu_m(w)^q.
\]
则有闭式展开
\[
H_q(\mu_m)
=
m\log\varphi+c(q)+O\!\Bigl(\Bigl(\frac{\varphi}{2}\Bigr)^m\Bigr),
\]
其中
\[
c(q):=\frac{1}{1-q}\log\!\Bigl(\frac{\varphi+\varphi^{-q}}{\sqrt5}\Bigr).
\]
并且 \(c(q)\) 在 \(q=1\) 处可连续延拓为
\[
c(1)=\frac{\log\varphi}{\varphi\sqrt5},
\]
相应地 Shannon 熵满足
\[
H(\mu_m)=m\log\varphi+\frac{\log\varphi}{\varphi\sqrt5}+o(1).
\]
因此 bin-fold 的全部 Rényi 熵族不仅具有统一的塌缩斜率 \(\log\varphi\)，而且其二阶常数项亦可闭式写出。
\end{theorem}
\begin{proof}
令 \(u_m(w):=1/|X_m|\) 为 \(X_m\) 上的均匀分布。
由定理 \ref{thm:fold-bin-renyi-divergence-limit}，对每个固定 \(q>0\)、\(q\neq 1\)，有
\[
D_q(\mu_m\Vert u_m)=D_q^\infty+O\!\Bigl(\Bigl(\frac{\varphi}{2}\Bigr)^m\Bigr),
\]
其中
\[
D_q^\infty
=
\frac{(2q-2)\log\varphi+\log(\varphi+\varphi^{-q})-q\log\sqrt5}{q-1}.
\]
由于 \(u_m\) 为均匀分布，Rényi 散度与 Rényi 熵满足恒等式
\[
D_q(\mu_m\Vert u_m)=\log|X_m|-H_q(\mu_m).
\]
又因 \(|X_m|=F_{m+2}\)，Binet 公式给出
\[
\log|X_m|=(m+2)\log\varphi-\frac12\log 5+O(\varphi^{-2m}).
\]
代入即得
\[
H_q(\mu_m)
=
m\log\varphi+
\Bigl(2\log\varphi-\frac12\log 5-D_q^\infty\Bigr)
+O\!\Bigl(\Bigl(\frac{\varphi}{2}\Bigr)^m\Bigr).
\]
把括号内常数化简，得到
\[
2\log\varphi-\frac12\log 5-D_q^\infty
=
\frac{1}{1-q}\log\!\Bigl(\frac{\varphi+\varphi^{-q}}{\sqrt5}\Bigr)
=c(q).
\]

当 \(q=1\) 时，命题 \ref{prop:fold-bin-kl-constant} 给出
\[
H(\mu_m)=\log|X_m|-D_{\mathrm{KL}}(\mu_m\Vert u_m)
=m\log\varphi+\frac{\log\varphi}{\varphi\sqrt5}+o(1).
\]
最后，对 \(c(q)\) 在 \(q=1\) 处应用 l'Hospital 法则，便得
\[
\lim_{q\to 1}c(q)=\frac{\log\varphi}{\varphi\sqrt5}.
\]
证毕。
\end{proof}

\paragraph{escort 温度轴上的末位 Gibbs 极限、两原子涨落与熵常数闭式}
前两条处理原始推前分布 \(\mu_m=\mu_{m,1}\)。
对整条 escort 温度轴，同一两态结构实际上同时控制末位 Gibbs 偏置、对数多重度的极限分布、一阶局部信息修正、Rényi--Shannon 常数以及局部指数谱。

\begin{theorem}[escort 族的末位 Gibbs 极限]\label{thm:xi-fold-escort-lastbit-gibbs-limit}
对任意固定 \(q\ge 0\)，定义 escort 测度
\[
\mu_{m,q}(w):=\frac{(d_m^{\mathrm{bin}}(w))^{q}}{S_q^{\mathrm{bin}}(m)},
\qquad
w\in X_m.
\]
则有极限
\[
\mu_{m,q}(w_m=1)\longrightarrow \frac{1}{1+\varphi^{q+1}},
\qquad
\mu_{m,q}(w_m=0)\longrightarrow \frac{\varphi^{q+1}}{1+\varphi^{q+1}}.
\]
特别地，\(\mu_{m,1}=\mu_m\)，并且
\[
\mu_m(w_m=1)\longrightarrow \frac{1}{1+\varphi^2}=\frac{1}{\varphi\sqrt5},
\qquad
\mu_m(w_m=0)\longrightarrow \frac{\varphi^2}{1+\varphi^2}=\frac{\varphi}{\sqrt5}.
\]
\end{theorem}
\begin{proof}
该极限正是命题 \ref{prop:fold-bin-escort-lastbit} 的第一部分。
对 \(q=1\)，只需使用恒等式 \(1+\varphi^2=\varphi\sqrt5\) 即可得到所示化简。
\end{proof}

\begin{theorem}[escort 对数多重度的两原子极限分布]\label{thm:xi-fold-escort-log-multiplicity-two-atom}
\leanverified{paper\_xi\_fold\_escort\_log\_multiplicity\_two\_atom}
沿用定理 \ref{thm:xi-fold-escort-lastbit-gibbs-limit} 的记号，并令 \(W_{m,q}\sim\mu_{m,q}\)。
定义中心化对数多重度
\[
X_m^{(q)}:=\log d_m^{\mathrm{bin}}(W_{m,q})-m\log\!\Bigl(\frac{2}{\varphi}\Bigr).
\]
则
\[
X_m^{(q)}\Longrightarrow \nu_q
:=
\frac{\varphi^{q+1}}{1+\varphi^{q+1}}\delta_{0}
+\frac{1}{1+\varphi^{q+1}}\delta_{-\log\varphi}.
\]
等价地，对任意固定 \(s\in\CC\)，有
\[
\lim_{m\to\infty}\mathbb E_{\mu_{m,q}}\!\bigl[e^{sX_m^{(q)}}\bigr]
=
\frac{\varphi^{q+1}+\varphi^{-s}}{1+\varphi^{q+1}}.
\]
从而
\[
\lim_{m\to\infty}\mathbb E_{\mu_{m,q}}[X_m^{(q)}]
=-\frac{\log\varphi}{1+\varphi^{q+1}},
\]
\[
\lim_{m\to\infty}\operatorname{Var}_{\mu_{m,q}}(X_m^{(q)})
=
\frac{\varphi^{q+1}}{(1+\varphi^{q+1})^2}(\log\varphi)^2.
\]
\end{theorem}
\begin{proof}
由定理 \ref{thm:fold-bin-two-state-asymptotic} 的一致对数展开，
\[
\log d_m^{\mathrm{bin}}(w)
=
m\log\!\Bigl(\frac{2}{\varphi}\Bigr)-w_m\log\varphi+O\!\Bigl(\Bigl(\frac{\varphi}{2}\Bigr)^m\Bigr)
\]
对 \(w\in X_m\) 一致成立，因此
\[
X_m^{(q)}=-W_{m,q}(m)\log\varphi+o_{\mathbb P}(1).
\]
再与定理 \ref{thm:xi-fold-escort-lastbit-gibbs-limit} 的末位极限分布合并，即得所示两原子极限。
变换公式、均值与方差都只是该两点分布的直接计算。
\end{proof}

\begin{corollary}[局部信息密度的一阶离散修正律]\label{cor:xi-fold-local-information-density-first-order-discrete-law}
令 \(W_m\sim\mu_m\)，并定义
\[
\gamma_m:=\frac{1}{m}\log d_m^{\mathrm{bin}}(W_m).
\]
再令 \(B\) 为参数
\[
\mathbb P(B=1)=\frac{1}{\varphi\sqrt5},
\qquad
\mathbb P(B=0)=\frac{\varphi}{\sqrt5}
\]
的 Bernoulli 变量。则有分布收敛
\[
m\left(\gamma_m-\log\!\frac{2}{\varphi}\right)
\Longrightarrow
-(\log\varphi)B.
\]
并且
\[
\lim_{m\to\infty}
m\left(\mathbb E[\gamma_m]-\log\!\frac{2}{\varphi}\right)
=
-\frac{\log\varphi}{\varphi\sqrt5},
\]
\[
\lim_{m\to\infty}
m^2\operatorname{Var}(\gamma_m)
=
(\log\varphi)^2
\frac{1}{\varphi\sqrt5}
\left(
1-\frac{1}{\varphi\sqrt5}
\right).
\]
因此 \(\gamma_m\) 的首个非平凡修正并不呈现 \(\sqrt m\)-Gaussian 涨落，而是在 \(m^{-1}\) 尺度上塌缩为末位几何所决定的两点律。
\end{corollary}
\begin{proof}
在定理 \ref{thm:xi-fold-escort-log-multiplicity-two-atom} 中取 \(q=1\)。
由于 \(\mu_{m,1}=\mu_m\)，且
\[
X_m^{(1)}
=
\log d_m^{\mathrm{bin}}(W_m)-m\log\!\Bigl(\frac{2}{\varphi}\Bigr)
=
m\left(\gamma_m-\log\!\frac{2}{\varphi}\right),
\]
故该定理立即给出
\[
m\left(\gamma_m-\log\!\frac{2}{\varphi}\right)
\Longrightarrow
\frac{\varphi}{\sqrt5}\delta_0+\frac{1}{\varphi\sqrt5}\delta_{-\log\varphi},
\]
这正是随机变量 \(-(\log\varphi)B\) 的分布。

又由定理 \ref{thm:xi-fold-escort-log-multiplicity-two-atom} 的证明，
\[
X_m^{(1)}=-W_m(m)\log\varphi+O\!\Bigl(\Bigl(\frac{\varphi}{2}\Bigr)^m\Bigr),
\]
且误差一致于样本点，从而 \((X_m^{(1)})_{m\ge 1}\) 一致有界。
因此其一、二阶矩都收敛到极限分布的相应矩。于是
\[
\lim_{m\to\infty}\mathbb E[X_m^{(1)}]
=
\mathbb E[-(\log\varphi)B]
=
-\frac{\log\varphi}{\varphi\sqrt5},
\]
\[
\lim_{m\to\infty}\operatorname{Var}(X_m^{(1)})
=
\operatorname{Var}((\log\varphi)B)
=
(\log\varphi)^2
\frac{1}{\varphi\sqrt5}
\left(
1-\frac{1}{\varphi\sqrt5}
\right).
\]
最后用
\[
X_m^{(1)}=m\left(\gamma_m-\log\!\frac{2}{\varphi}\right)
\]
即可得到所示均值与方差极限。证毕。
\end{proof}


\begin{theorem}[escort 温度族的 Rényi--Shannon 常数闭式]\label{thm:xi-fold-escort-renyi-shannon-constants}
固定 \(q\ge 0\)，并对任意 \(r>0\)、\(r\neq 1\) 定义 Rényi 熵
\[
H_r(\mu_{m,q}):=\frac{1}{1-r}\log\sum_{w\in X_m}\mu_{m,q}(w)^r.
\]
则有展开
\[
H_r(\mu_{m,q})=m\log\varphi+C_{q,r}+O\!\Bigl(\Bigl(\frac{\varphi}{2}\Bigr)^m\Bigr),
\]
其中
\[
C_{q,r}
:=
\frac{1}{1-r}\left[
\log\!\Bigl(\frac{\varphi+\varphi^{-qr}}{\sqrt5}\Bigr)
-r\log\!\Bigl(\frac{\varphi+\varphi^{-q}}{\sqrt5}\Bigr)
\right].
\]
相应的 Shannon 熵满足
\[
H(\mu_{m,q})=m\log\varphi+C_q+O\!\Bigl(\Bigl(\frac{\varphi}{2}\Bigr)^m\Bigr),
\]
其中
\[
C_q
:=
\log\!\Bigl(\frac{\varphi+\varphi^{-q}}{\sqrt5}\Bigr)
+\frac{q\log\varphi}{1+\varphi^{q+1}}.
\]
\end{theorem}
\begin{proof}
由定义
\[
\sum_{w\in X_m}\mu_{m,q}(w)^r=\frac{S_{qr}^{\mathrm{bin}}(m)}{(S_q^{\mathrm{bin}}(m))^r}.
\]
记
\[
A(a):=\log\!\Bigl(\frac{\varphi+\varphi^{-a}}{\sqrt5}\Bigr).
\]
由命题 \ref{prop:fold-bin-power-sum-asymptotic}，
\[
\log S_a^{\mathrm{bin}}(m)
=
am\log\!\Bigl(\frac{2}{\varphi}\Bigr)+m\log\varphi+A(a)
+O\!\Bigl(\Bigl(\frac{\varphi}{2}\Bigr)^m\Bigr).
\]
因此
\[
H_r(\mu_{m,q})
=
m\log\varphi+\frac{A(qr)-rA(q)}{1-r}
+O\!\Bigl(\Bigl(\frac{\varphi}{2}\Bigr)^m\Bigr),
\]
即得 \(C_{q,r}\)。

对 Shannon 熵，利用恒等式
\[
H(\mu_{m,q})=\log S_q^{\mathrm{bin}}(m)-q\,\mathbb E_{\mu_{m,q}}\!\bigl[\log d_m^{\mathrm{bin}}(W_{m,q})\bigr]
\]
并代入命题 \ref{prop:fold-bin-escort-lastbit} 给出的纤维信息展开，即得所示 \(C_q\)。
\end{proof}

\begin{theorem}[escort 全族的严格单点多重分形塌缩]\label{thm:xi-fold-escort-single-point-multifractal-collapse}
对每个固定 \(q\ge 0\)，定义局部指数
\[
\alpha_{m,q}(w):=-\frac{1}{m}\log\mu_{m,q}(w),
\qquad
w\in X_m.
\]
则有一致极限
\[
\sup_{w\in X_m}\bigl|\alpha_{m,q}(w)-\log\varphi\bigr|\longrightarrow 0.
\]
更精确地，
\[
\alpha_{m,q}(w)
=
\log\varphi
+\frac{q\,w_m\log\varphi+\log\!\bigl(\frac{\varphi+\varphi^{-q}}{\sqrt5}\bigr)}{m}
+o\!\Bigl(\frac{1}{m}\Bigr)
\]
在 \(w\in X_m\) 上一致成立。

若进一步定义离散谱函数
\[
f_q(\alpha)
:=
\lim_{\varepsilon\downarrow 0}\limsup_{m\to\infty}
\frac{1}{m}\log
\#\bigl\{w\in X_m:\ |\alpha_{m,q}(w)-\alpha|<\varepsilon\bigr\},
\]
则
\[
f_q(\alpha)=
\begin{cases}
\log\varphi,& \alpha=\log\varphi,\\[1mm]
-\infty,& \alpha\neq \log\varphi.
\end{cases}
\]
\end{theorem}
\begin{proof}
由定义
\[
\alpha_{m,q}(w)=-\frac{q}{m}\log d_m^{\mathrm{bin}}(w)+\frac{1}{m}\log S_q^{\mathrm{bin}}(m).
\]
再将定理 \ref{thm:fold-bin-two-state-asymptotic} 的一致对数展开与上一个证明中的
\(\log S_q^{\mathrm{bin}}(m)\) 展开代入，即得所示统一公式。
因此 \(\sup_{w}|\alpha_{m,q}(w)-\log\varphi|\to 0\)。

若 \(\alpha\neq\log\varphi\)，取 \(\varepsilon<|\alpha-\log\varphi|/2\)，则对充分大 \(m\) 上述集合为空，
故 \(f_q(\alpha)=-\infty\)。
若 \(\alpha=\log\varphi\)，则对任意固定 \(\varepsilon>0\)，充分大 \(m\) 时全部 \(w\in X_m\) 都落入该带，于是
\[
f_q(\log\varphi)=\lim_{m\to\infty}\frac{1}{m}\log|X_m|=\log\varphi.
\]
证毕。
\end{proof}

\begin{corollary}[escort--uniform 相对熵的显式温度律]\label{cor:xi-fold-escort-uniform-kl-temperature-law}
令 \(u_m\) 为 \(X_m\) 上的均匀分布。则对每个固定 \(q\ge 0\)，有
\[
D_{\mathrm{KL}}(\mu_{m,q}\Vert u_m)
=
D(q)+O\!\Bigl(\Bigl(\frac{\varphi}{2}\Bigr)^m\Bigr),
\]
其中
\[
D(q)
:=
2\log\varphi-\log(\varphi+\varphi^{-q})
-\frac{q\log\varphi}{1+\varphi^{q+1}}.
\]
特别地，当 \(q=1\) 时，
\[
D(1)=2\log\varphi-\frac12\log 5-\frac{\log\varphi}{\varphi\sqrt5},
\]
从而恢复命题 \ref{prop:fold-bin-kl-constant} 的极限常数。
\end{corollary}
\begin{proof}
由恒等式
\[
D_{\mathrm{KL}}(\mu_{m,q}\Vert u_m)=\log|X_m|-H(\mu_{m,q})
\]
以及定理 \ref{thm:xi-fold-escort-renyi-shannon-constants}，
\[
\log|X_m|=(m+2)\log\varphi-\frac12\log 5+O(\varphi^{-2m}),
\]
即可得到所示公式。
当 \(q=1\) 时，只需再用 \(\varphi+\varphi^{-1}=\sqrt5\) 与 \(1+\varphi^2=\varphi\sqrt5\) 化简。
\end{proof}

\begin{theorem}[有限素数局部化的 \(p\)-进核刚性分类]\label{thm:xi-localized-integers-padic-kernel-rigidity}
对有限素数集 \(S\subset\mathbb P\)，记
\[
G_S:=\ZZ[S^{-1}],
\qquad
K_S:=\prod_{p\in S}\ZZ_p.
\]
则有短正合列
\[
0\longrightarrow K_S\longrightarrow \widehat{G_S}\longrightarrow \TT\longrightarrow 0,
\qquad
(\widehat{G_S})^0\cong \TT,
\qquad
\cdim_{\mathrm{fr}}(G_S)=1.
\]
并且对任意有限素数集 \(S,T\subset\mathbb P\)，成立
\[
G_S\cong G_T
\iff
K_S\cong K_T
\iff
S=T.
\]
因此有限局部化群的连通相位部分完全相同，而其全部可分辨算术信息都由全不连通
\(p\)-进核 \(K_S\) 完全承载。
\end{theorem}
\begin{proof}
定理 \ref{thm:cdim-star-z1s-dual-extension} 给出短正合列
\[
0\longrightarrow K_S\longrightarrow \widehat{G_S}\longrightarrow \TT\longrightarrow 0.
\]
由于 \(K_S\) 全不连通，连通分量只能映到 \(\TT\) 的连通分量，故
\(
(\widehat{G_S})^0\cong \TT
\)。
另一方面，\(G_S\otimes_{\ZZ}\QQ\cong \QQ\) 为秩 \(1\) 的有限秩无挠群，故由定理
\ref{thm:conclusion-cdim-finite-rank-extension}，
\[
\cdim_{\mathrm{fr}}(G_S)=1.
\]

若 \(K_S\cong K_T\)，则对任意素数 \(p\) 考察商群 \(K_S/pK_S\)。
当 \(p\in S\) 时，
\[
\ZZ_p/p\ZZ_p\cong \FF_p\neq 0;
\]
而当 \(q\neq p\) 时，\(p\) 在 \(\ZZ_q\) 中可逆，故
\[
\ZZ_q/p\ZZ_q=0.
\]
于是
\[
K_S/pK_S\neq 0
\iff
p\in S.
\]
同理
\(
K_T/pK_T\neq 0\iff p\in T
\)。
由 \(K_S\cong K_T\) 可知上述商群同时非零或同时为零，故 \(S=T\)。

反向蕴含 \(S=T\Rightarrow G_S\cong G_T\) 与 \(S=T\Rightarrow K_S\cong K_T\) 显然。
而 \(G_S\cong G_T\Rightarrow S=T\) 亦可由推论
\ref{cor:xi-localized-integers-order-dual-quotient} 直接得到。
故三者等价。证毕。
\end{proof}
\leanverified{paper\_xi\_localized\_integers\_padic\_kernel\_rigidity}

\begin{corollary}[有限局部化塔的逐素数增量律]\label{cor:xi-localized-integers-prime-axis-increment-law}
设 \(S\subset T\subset\mathbb P\) 为有限素数集。则有自然短正合列
\[
0\longrightarrow G_S\longrightarrow G_T\longrightarrow
\bigoplus_{p\in T\setminus S}\ZZ(p^\infty)\longrightarrow 0,
\]
其 Pontryagin 对偶为
\[
0\longrightarrow \prod_{p\in T\setminus S}\ZZ_p
\longrightarrow \widehat{G_T}
\longrightarrow \widehat{G_S}
\longrightarrow 0.
\]
特别地，每加入一个新素数 \(p\)，对偶端恰增加一条新的 \(p\)-进轴，而连通商
\(\TT\) 与有限秩圆维 \(\cdim_{\mathrm{fr}}=1\) 保持不变。
\end{corollary}
\begin{proof}
由
\[
0\to \ZZ\to G_T\to \bigoplus_{p\in T}\ZZ(p^\infty)\to 0,
\qquad
0\to \ZZ\to G_S\to \bigoplus_{p\in S}\ZZ(p^\infty)\to 0
\]
及包含 \(G_S\subset G_T\) 可得
\[
G_T/G_S
\cong
(G_T/\ZZ)/(G_S/\ZZ)
\cong
\bigoplus_{p\in T\setminus S}\ZZ(p^\infty),
\]
从而得到第一条短正合列。对偶化并使用
\(
\widehat{\ZZ(p^\infty)}\cong \ZZ_p
\)
即得第二条。

最后，由定理 \ref{thm:xi-localized-integers-padic-kernel-rigidity}，
\(\widehat{G_S}\) 与 \(\widehat{G_T}\) 的连通分量均同构于 \(\TT\)，且
\(\cdim_{\mathrm{fr}}(G_S)=\cdim_{\mathrm{fr}}(G_T)=1\)。
故新增信息完全进入全不连通核。证毕。
\end{proof}
\leanverified{paper\_xi\_localized\_integers\_prime\_axis\_increment\_law}

\begin{theorem}[连通与有理不变量的统一不可区分性]\label{thm:xi-localized-integers-connected-rational-blindness}
设 \(\mathfrak I\) 为定义在有限素数局部化群族
\[
\{G_S=\ZZ[S^{-1}]:\ S\subset\mathbb P,\ |S|<\infty\}
\]
上的同构不变量，并满足下列任一条件：
\begin{enumerate}
  \item \(\mathfrak I(G)\) 仅依赖于 \(G\otimes_{\ZZ}\QQ\)；
  \item \(\mathfrak I(G)\) 仅依赖于 \((\widehat G)^0\)；
  \item \(\mathfrak I(G)\) 仅依赖于 \(\cdim_{\mathrm{fr}}(G)\)。
\end{enumerate}
则对任意有限 \(S,T\subset\mathbb P\) 都有
\[
\mathfrak I(G_S)=\mathfrak I(G_T).
\]
特别地，任何只读取有理部分、连通部分或有限秩圆维的一维不变量，都不可能对该族给出完备分类。
\end{theorem}
\begin{proof}
对任意有限 \(S\subset\mathbb P\)，都有
\[
G_S\otimes_{\ZZ}\QQ\cong \QQ.
\]
又由定理 \ref{thm:xi-localized-integers-padic-kernel-rigidity}，
\[
(\widehat{G_S})^0\cong \TT,
\qquad
\cdim_{\mathrm{fr}}(G_S)=1.
\]
因此若 \(\mathfrak I\) 只通过上述任一对象因子化，则它在所有 \(G_S\) 上都取同一数值。

另一方面，定理 \ref{thm:xi-localized-integers-padic-kernel-rigidity} 已证明
\[
G_S\cong G_T
\iff
S=T.
\]
故这类不变量都不能区分不同的素数支撑，更不可能给出完备分类。证毕。
\end{proof}
\leanverified{paper\_xi\_localized\_integers\_connected\_rational\_blindness}

\paragraph{纤维群胚代数的中心谱、Fuglede--Kadison 体积、容量截断与 window-\texorpdfstring{$6$}{6} 异常分层}
下列结果把命题 \ref{prop:fold-bin-groupoid-algebra-wedderburn}、定理
\ref{thm:op-algebra-fold-watatani-index-equals-multiplicity-field}、定理
\ref{thm:fold-oracle-capacity-closed-form}、命题
\ref{prop:fold-oracle-capacity-bimeasure-identity}、定理
\ref{thm:fold-bin-gauge-volume-stirling}、定理
\ref{thm:xi-foldbin-gauge-representation-distribution-law-m6-spectrum} 与定理
\ref{thm:xi-window6-center-exact-splitting-bdry-cyclic} 组织为同一条中心算子链：纤维多重度场既是群胚代数中心条件期望的最小中心值指数元，也是规范体积、离散容量与
window-\(6\) 异常签名分层的统一谱变量。

\begin{theorem}[纤维群胚代数的最小中心值指数元与谱顶端]\label{thm:xi-foldbin-minimal-center-valued-index-spectrum-top}
设
\[
A_m:=\CC[\mathcal R_m^{\mathrm{bin}}]
\cong
\bigoplus_{w\in X_m}M_{d_m^{\mathrm{bin}}(w)}(\CC),
\qquad
Z_m:=Z(A_m)=\bigoplus_{w\in X_m}\CC z_w,
\]
其中 \(z_w\) 为对应于第 \(w\)-块的最小中心幂等元。定义逐块归一化迹条件期望
\[
E_m\!\left(\bigoplus_{w\in X_m}a_w\right)
:=
\bigoplus_{w\in X_m}
\frac{\operatorname{Tr}(a_w)}{d_m^{\mathrm{bin}}(w)}I_{d_m^{\mathrm{bin}}(w)},
\]
并记
\[
d_w:=d_m^{\mathrm{bin}}(w),
\qquad
\Lambda_m:=\sum_{w\in X_m}d_w z_w\in Z_m.
\]
再定义
\[
M_m:=\max_{w\in X_m}d_w,
\qquad
K_m:=\#\{w\in X_m:\ d_w=M_m\}.
\]
则有：
\begin{enumerate}
  \item 对任意 \(x\in A_m^+\)，都有
  \[
  E_m(x)\ge \Lambda_m^{-1}x,
  \qquad
  \Lambda_m^{-1}=\sum_{w\in X_m}d_w^{-1}z_w.
  \]
  \item 若 \(\Gamma=\sum_{w\in X_m}\gamma_w z_w\in Z_m\) 满足 \(\gamma_w>0\) 且
  \[
  E_m(x)\ge \Gamma^{-1}x
  \qquad
  \forall x\in A_m^+,
  \]
  则 \(\Gamma\ge \Lambda_m\)。因而 \(\Lambda_m\) 是 \(E_m\) 的最小中心值指数元。
  \item \(\Lambda_m\) 的标量指数恰为
  \[
  \|\Lambda_m\|_\infty=M_m,
  \]
  且中心左乘算子
  \[
  L_{\Lambda_m}:Z_m\to Z_m,
  \qquad
  z\mapsto \Lambda_m z
  \]
  在特征值 \(M_m\) 处的几何重数恰为 \(K_m\)。
  \item 若极限
  \[
  \alpha_*:=\lim_{m\to\infty}\frac{1}{m}\log M_m,
  \qquad
  g_*:=\lim_{m\to\infty}\frac{1}{m}\log K_m
  \]
  存在，则
  \[
  \alpha_*=
  \lim_{m\to\infty}\frac{1}{m}\log\|\Lambda_m\|_\infty,
  \qquad
  g_*=
  \lim_{m\to\infty}\frac{1}{m}\log
  \dim\ker\!\bigl(M_m I_{Z_m}-L_{\Lambda_m}\bigr).
  \]
\end{enumerate}
\end{theorem}
\begin{proof}
对任意 \(x\in A_m^+\) 写成块对角形式
\[
x=\bigoplus_{w\in X_m}x_w,
\qquad
x_w\in M_{d_w}(\CC)^+.
\]
在单个矩阵块上，
\[
E_m(x)_w=\frac{\operatorname{Tr}(x_w)}{d_w}I_{d_w}.
\]
又对任意正矩阵 \(x_w\)，都有基本不等式
\[
x_w\le \operatorname{Tr}(x_w)\,I_{d_w}.
\]
因此
\[
E_m(x)_w\ge \frac{1}{d_w}x_w.
\]
逐块并合即得
\[
E_m(x)\ge \left(\sum_{w\in X_m}d_w^{-1}z_w\right)x=\Lambda_m^{-1}x,
\]
从而得到第 \(1\) 条。

现证最小性。设 \(\Gamma=\sum_w\gamma_w z_w\) 满足 \(\gamma_w>0\) 且
\[
E_m(x)\ge \Gamma^{-1}x
\qquad
\forall x\in A_m^+.
\]
固定 \(w\in X_m\)，取第 \(w\)-块中的秩一投影 \(p_w\)。则
\[
E_m(p_w)=\frac{1}{d_w}I_{d_w}.
\]
若 \(\gamma_w<d_w\)，则在 \(\operatorname{Ran}(p_w)\) 上有
\[
\left(E_m(p_w)-\Gamma^{-1}p_w\right)\big|_{\operatorname{Ran}(p_w)}
=
\left(\frac{1}{d_w}-\frac{1}{\gamma_w}\right)I_{\operatorname{Ran}(p_w)},
\]
该算子严格为负，矛盾。故 \(\gamma_w\ge d_w\) 对所有 \(w\) 都成立，即
\(\Gamma\ge \Lambda_m\)。第 \(2\) 条得证。

第 \(3\) 条是 \(\Lambda_m\) 在中心基 \(\{z_w\}_{w\in X_m}\) 上按坐标对角的直接结果：其特征值恰为 \(\{d_w\}_{w\in X_m}\)，故算子范数等于最大坐标 \(M_m\)，而顶端特征空间由满足 \(d_w=M_m\) 的坐标轴张成，维数正是 \(K_m\)。

第 \(4\) 条只是把第 \(3\) 条的两条精确恒等式写成指数增长率。证毕。
\end{proof}

\begin{theorem}[中心值指数元的 Fuglede--Kadison 对数体积与规范体积常数差]\label{thm:xi-foldbin-fkdet-logvolume-kappa}
沿用定理 \ref{thm:xi-foldbin-minimal-center-valued-index-spectrum-top} 的记号，并记
\[
\mathcal G_m^{\mathrm{bin}}:=\Aut(\Fold_m^{\mathrm{bin}}),
\qquad
g_m:=2^{-m}\log|\mathcal G_m^{\mathrm{bin}}|,
\]
\[
\kappa_m:=2^{-m}\sum_{w\in X_m}d_w\log d_w.
\]
在 \(A_m\) 上定义概率迹
\[
\tau_m\!\left(\bigoplus_{w\in X_m}a_w\right)
:=
2^{-m}\sum_{w\in X_m}\operatorname{Tr}(a_w).
\]
则
\[
\log\Delta_{\tau_m}(\Lambda_m):=\tau_m(\log\Lambda_m)=\kappa_m.
\]
因此
\[
|\mathcal G_m^{\mathrm{bin}}|^{1/2^m}
=
e^{-1+O\!\left(m\left(\frac{\varphi}{2}\right)^m\right)}
\Delta_{\tau_m}(\Lambda_m),
\]
并且
\[
\lim_{m\to\infty}
\frac{|\mathcal G_m^{\mathrm{bin}}|^{1/2^m}}{\Delta_{\tau_m}(\Lambda_m)}
=
e^{-1}.
\]
\end{theorem}
\begin{proof}
由 \(\sum_{w\in X_m}d_w=2^m\)，有
\[
\tau_m(1_{A_m})=2^{-m}\sum_{w\in X_m}d_w=1,
\]
故 \(\tau_m\) 确为概率迹。又因 \(\Lambda_m\) 在第 \(w\)-块上等于 \(d_w I_{d_w}\)，故其函数演算满足
\[
\log\Lambda_m=\sum_{w\in X_m}(\log d_w)\,z_w.
\]
于是
\[
\tau_m(\log\Lambda_m)
=
2^{-m}\sum_{w\in X_m}\operatorname{Tr}\!\bigl((\log d_w)I_{d_w}\bigr)
=
2^{-m}\sum_{w\in X_m}d_w\log d_w
=
\kappa_m,
\]
从而
\[
\log\Delta_{\tau_m}(\Lambda_m)=\kappa_m.
\]

另一方面，定理 \ref{thm:fold-bin-gauge-volume-stirling} 给出
\[
g_m=\kappa_m-1+O\!\left(m\left(\frac{\varphi}{2}\right)^m\right).
\]
对两边取指数，并用
\[
e^{g_m}=|\mathcal G_m^{\mathrm{bin}}|^{1/2^m},
\qquad
e^{\kappa_m}=\Delta_{\tau_m}(\Lambda_m),
\]
即得
\[
|\mathcal G_m^{\mathrm{bin}}|^{1/2^m}
=
e^{-1+O\!\left(m\left(\frac{\varphi}{2}\right)^m\right)}
\Delta_{\tau_m}(\Lambda_m).
\]
令 \(m\to\infty\) 即得极限比值 \(e^{-1}\)。证毕。
\end{proof}

\begin{theorem}[预言机容量是中心值指数元的谱截断迹]\label{thm:xi-foldbin-oracle-capacity-central-spectral-truncation}
沿用定理 \ref{thm:xi-foldbin-minimal-center-valued-index-spectrum-top} 的记号。对整数 \(B\ge 0\)，定义中心谱截断
\[
\Lambda_m\wedge 2^B
:=
\sum_{w\in X_m}\min(d_w,2^B)\,z_w\in Z_m,
\]
并记
\[
1_{\Lambda_m\le 2^B}:=\sum_{d_w\le 2^B}z_w,
\qquad
1_{\Lambda_m>2^B}:=\sum_{d_w>2^B}z_w.
\]
令 \(C_m(B)\) 为定理 \ref{thm:fold-oracle-capacity-closed-form} 的离散预言机容量。再定义中心上的三个迹
\[
\operatorname{Tr}_{Z_m}\!\left(\sum_{w\in X_m}c_w z_w\right):=\sum_{w\in X_m}c_w,
\qquad
u_m^Z:=|X_m|^{-1}\operatorname{Tr}_{Z_m},
\]
\[
\tau_m^Z\!\left(\sum_{w\in X_m}c_w z_w\right)
:=
2^{-m}\sum_{w\in X_m}d_w c_w.
\]
则有严格恒等式
\[
C_m(B)=\operatorname{Tr}_{Z_m}\!\bigl(\Lambda_m\wedge 2^B\bigr),
\]
以及
\[
\frac{C_m(B)}{2^m}
=
\tau_m^Z\!\bigl(1_{\Lambda_m\le 2^B}\bigr)
+
2^{B-m}|X_m|\,u_m^Z\!\bigl(1_{\Lambda_m>2^B}\bigr).
\]
特别地，
\[
C_m(B)\le 2^B\dim Z_m=2^B|X_m|.
\]
\end{theorem}
\begin{proof}
由于 \(\Lambda_m\) 在中心基 \(\{z_w\}\) 上对角，且其谱值恰为 \(\{d_w\}_{w\in X_m}\)，函数演算直接给出
\[
\Lambda_m\wedge 2^B
=
\sum_{w\in X_m}\min(d_w,2^B)\,z_w.
\]
对其取计数迹 \(\operatorname{Tr}_{Z_m}\) 得
\[
\operatorname{Tr}_{Z_m}\!\bigl(\Lambda_m\wedge 2^B\bigr)
=
\sum_{w\in X_m}\min(d_w,2^B).
\]
而右端正是定理 \ref{thm:fold-oracle-capacity-closed-form} 的容量闭式，故得到第一条恒等式。

再将上式按 \(d_w\le 2^B\) 与 \(d_w>2^B\) 分拆并除以 \(2^m\)，得
\[
\frac{C_m(B)}{2^m}
=
2^{-m}\sum_{d_w\le 2^B}d_w
+
2^{-m}\sum_{d_w>2^B}2^B.
\]
第一项恰为
\[
\tau_m^Z\!\bigl(1_{\Lambda_m\le 2^B}\bigr),
\]
第二项则等于
\[
2^{B-m}\#\{w\in X_m:\ d_w>2^B\}
=
2^{B-m}|X_m|\,u_m^Z\!\bigl(1_{\Lambda_m>2^B}\bigr),
\]
于是得到第二条公式。这正是命题 \ref{prop:fold-oracle-capacity-bimeasure-identity} 在中心谱上的算子化改写。

最后，由点态不等式
\[
\Lambda_m\wedge 2^B\le 2^B\,1_{Z_m}
\]
取 \(\operatorname{Tr}_{Z_m}\) 即得
\[
C_m(B)\le 2^B\operatorname{Tr}_{Z_m}(1_{Z_m})=2^B|X_m|.
\]
证毕。
\end{proof}

\begin{theorem}[window-\texorpdfstring{$6$}{6} 异常签名空间的规范三层分解]\label{thm:xi-window6-anomaly-signature-three-layer-splitting}
沿用定理 \ref{thm:xi-window6-center-exact-splitting-bdry-cyclic}、定理
\ref{thm:xi-foldbin-gauge-representation-distribution-law-m6-spectrum} 与定义
\ref{def:fold-bin-gauge-sign} 的记号。记
\[
C_{\partial}:=(\ZZ/2\ZZ)^{X_6^{\mathrm{bdry}}}\cong(\ZZ/2\ZZ)^3,
\qquad
C_{\mathrm{int}}:=(\ZZ/2\ZZ)^{X_{6,\mathrm{cyc},2}}\cong(\ZZ/2\ZZ)^5,
\]
并定义
\[
C_{\mathrm{bulk}}
:=
\bigl(S_3^{\,4}\times S_4^{\,9}\bigr)^{\mathrm{ab}}
\cong
(\ZZ/2\ZZ)^4\oplus(\ZZ/2\ZZ)^9
\cong
(\ZZ/2\ZZ)^{13}.
\]
则
\[
Z\!\bigl(\mathcal G_6^{\mathrm{bin}}\bigr)
\cong
C_{\partial}\oplus C_{\mathrm{int}}
\cong
(\ZZ/2\ZZ)^3\oplus(\ZZ/2\ZZ)^5
\cong
(\ZZ/2\ZZ)^8.
\]
并且在符号映射 \(\mathrm{sgn}_6\) 给出的 \(21\) 比特坐标下，
\[
\bigl(\mathcal G_6^{\mathrm{bin}}\bigr)^{\mathrm{ab}}
\cong
C_{\partial}\oplus C_{\mathrm{int}}\oplus C_{\mathrm{bulk}}
\cong
(\ZZ/2\ZZ)^3\oplus(\ZZ/2\ZZ)^5\oplus(\ZZ/2\ZZ)^{13}.
\]
因此 window-\(6\) 的异常签名空间具有规范三层分解：\(3\) 维 boundary 中心层、\(5\) 维非 boundary 的
\(d=2\) 中心层，以及 \(13\) 维非中心 bulk 符号层。
\end{theorem}
\begin{proof}
第一条分解正是定理 \ref{thm:xi-window6-center-exact-splitting-bdry-cyclic} 的结论：全部中心恰由八条
\(d=2\) 纤维给出，并按 boundary 三条与其余五条非 boundary 的 \(d=2\) 纤维裂解为
\[
Z\!\bigl(\mathcal G_6^{\mathrm{bin}}\bigr)
\cong
C_{\partial}\oplus C_{\mathrm{int}}
\cong
(\ZZ/2\ZZ)^3\oplus(\ZZ/2\ZZ)^5.
\]

另一方面，由定理 \ref{thm:xi-foldbin-gauge-representation-distribution-law-m6-spectrum}，
\[
\mathcal G_6^{\mathrm{bin}}
\cong
S_2^{\,8}\times S_3^{\,4}\times S_4^{\,9},
\qquad
\bigl(\mathcal G_6^{\mathrm{bin}}\bigr)^{\mathrm{ab}}
\cong
(\ZZ/2\ZZ)^8\oplus(\ZZ/2\ZZ)^4\oplus(\ZZ/2\ZZ)^9.
\]
其中第一个 \((\ZZ/2\ZZ)^8\) 正对应全部 \(d=2\) 纤维的符号坐标；由上一步，它又规范地分裂为
\[
(\ZZ/2\ZZ)^8\cong C_{\partial}\oplus C_{\mathrm{int}}.
\]
其余 \((\ZZ/2\ZZ)^4\oplus(\ZZ/2\ZZ)^9\) 则来自 \(S_3^{\,4}\times S_4^{\,9}\) 的符号荷，构成
\(C_{\mathrm{bulk}}\)。

最后，由于 \(S_3\) 与 \(S_4\) 的中心都平凡，这 \(13\) 个比特不可能来自中心层；它们只在交换子商与符号映射口径下可见。故得到所述三层分解。证毕。
\end{proof}

\begin{theorem}[window-\texorpdfstring{$6$}{6} boundary 奇偶子格在交换化中的规范直和实现]\label{thm:xi-window6-boundary-parity-canonical-direct-summand}
沿用定理 \ref{thm:xi-window6-anomaly-signature-three-layer-splitting} 与定理
\ref{thm:xi-window6-center-exact-splitting-bdry-cyclic} 的记号，并记
\[
G_6^{\mathrm{bin}}:=\mathcal G_6^{\mathrm{bin}},
\qquad
Z_{\partial}:=\langle \tau_w:\ w\in X_6^{\mathrm{bdry}}\rangle
\cong
(\ZZ/2\ZZ)^3
\subset Z(G_6^{\mathrm{bin}}).
\]
若
\[
\operatorname{ab}_6:G_6^{\mathrm{bin}}\longrightarrow \bigl(G_6^{\mathrm{bin}}\bigr)^{\mathrm{ab}}
\]
为交换化映射，则 \(\operatorname{ab}_6|_{Z_{\partial}}\) 为单同态，其像恰为
\[
C_{\partial}=(\ZZ/2\ZZ)^{X_6^{\mathrm{bdry}}}.
\]
从而有规范直和分解
\[
\boxed{\;
\bigl(G_6^{\mathrm{bin}}\bigr)^{\mathrm{ab}}
\cong
C_{\partial}\oplus (\ZZ/2\ZZ)^{X_6\setminus X_6^{\mathrm{bdry}}}
\cong
(\ZZ/2\ZZ)^3\oplus(\ZZ/2\ZZ)^{18}.\;}
\]
其中第一因子由中心子群 \(Z_{\partial}\) 在交换化中的像规范实现。
\end{theorem}
\begin{proof}
定理 \ref{thm:xi-window6-center-exact-splitting-bdry-cyclic} 已将 boundary 三条 \(d=2\) 纤维所生成的中心层识别为
\[
C_{\partial}\cong (\ZZ/2\ZZ)^{X_6^{\mathrm{bdry}}}.
\]
因此三条中心 involution \(\tau_w\) 在交换化中的像彼此独立，\(\operatorname{ab}_6|_{Z_{\partial}}\) 为单同态，且其像正是 \(C_{\partial}\)。

另一方面，定理 \ref{thm:xi-window6-anomaly-signature-three-layer-splitting} 给出
\[
\bigl(G_6^{\mathrm{bin}}\bigr)^{\mathrm{ab}}
\cong
C_{\partial}\oplus C_{\mathrm{int}}\oplus C_{\mathrm{bulk}}
\cong
(\ZZ/2\ZZ)^3\oplus(\ZZ/2\ZZ)^5\oplus(\ZZ/2\ZZ)^{13}.
\]
将后两项合并即得
\[
C_{\mathrm{int}}\oplus C_{\mathrm{bulk}}
\cong
(\ZZ/2\ZZ)^{18}
\cong
(\ZZ/2\ZZ)^{X_6\setminus X_6^{\mathrm{bdry}}}.
\]
于是得到所示规范直和分解，其中第一因子恰由 \(Z_{\partial}\) 的像实现。证毕。
\end{proof}

\begin{corollary}[window-\texorpdfstring{$6$}{6} 的 boundary parity 是 \texorpdfstring{$21$}{21} 维异常签名空间的唯一三维最小忠实商]\label{cor:xi-window6-boundary-parity-unique-minimal-faithful-quotient}
沿用定理 \ref{thm:xi-window6-boundary-parity-canonical-direct-summand} 的记号，并记
\[
\mathcal A_6:=\bigl(G_6^{\mathrm{bin}}\bigr)^{\mathrm{ab}}
\cong
(\ZZ/2\ZZ)^{21},
\qquad
P_6:=C_{\partial}\cong(\ZZ/2\ZZ)^3.
\]
设 \(B\) 为任意初等 \(2\)-群，\(\phi:\mathcal A_6\to B\) 为群同态，并假设 \(\phi|_{P_6}\) 为单同态。则有：
\begin{enumerate}
  \item \(\dim_{\FF_2}B\ge 3\)；
  \item 当且仅当 \(\dim_{\FF_2}B=3\) 时，\(\phi\) 与沿某个 \(18\) 维补空间的投影等价。
\end{enumerate}
等价地，在保持全部 boundary parity 信息的前提下，\(\mathcal A_6\) 的最小忠实商恰为
\[
\mathcal A_6/K\cong P_6\cong (\ZZ/2\ZZ)^3,
\]
其中 \(K\subset \mathcal A_6\) 为某个 \(18\) 维补空间；因而其核维数必且仅必为 \(18\)。
\end{corollary}
\begin{proof}
由于 \(P_6\cong(\ZZ/2\ZZ)^3\) 且 \(\phi|_{P_6}\) 单射，像 \(\phi(P_6)\) 是 \(B\) 的一个三维 \(\FF_2\)-子空间，故必有 \(\dim_{\FF_2}B\ge 3\)。

若 \(\dim_{\FF_2}B=3\)，则 \(\phi|_{P_6}:P_6\to B\) 为同构。此时 \(\phi\) 的像即为 \(B\)，故 \(\operatorname{rank}\phi=3\)，从而
\[
\dim_{\FF_2}\ker\phi=21-3=18.
\]
又因 \(\phi|_{P_6}\) 单射，故
\[
\ker\phi\cap P_6=\{0\}.
\]
因此 \(\ker\phi\) 给出 \(P_6\) 的一个 \(18\) 维补空间，并有
\[
\mathcal A_6=P_6\oplus\ker\phi.
\]
于是 \(\phi\) 与沿 \(\ker\phi\) 的投影
\[
\mathcal A_6\twoheadrightarrow \mathcal A_6/\ker\phi\cong P_6
\]
在识别 \(P_6\cong B\) 后等价。

反之，若 \(K\subset \mathcal A_6\) 为任一与 \(P_6\) 互补的 \(18\) 维子空间，则自然投影
\[
\mathcal A_6\twoheadrightarrow \mathcal A_6/K\cong P_6\cong(\ZZ/2\ZZ)^3
\]
在 \(P_6\) 上显然保持忠实，且目标维数恰为 \(3\)。故三维正是最小可达维数。证毕。
\end{proof}
\leanverified{paper\_xi\_window6\_boundary\_parity\_unique\_minimal\_faithful\_quotient}

由此，纤维多重度场在群胚代数层面实现如下统一链：
\[
\CC[\mathcal R_m^{\mathrm{bin}}]
\Longrightarrow
\Lambda_m=\sum_{w\in X_m}d_w z_w
\Longrightarrow
\begin{cases}
\text{最小中心值指数元},\\
\log\Delta_{\tau_m}(\Lambda_m)=\kappa_m,\\
C_m(B)=\operatorname{Tr}_{Z_m}(\Lambda_m\wedge 2^B).
\end{cases}
\]
而在 \(m=6\) 时，\(\Lambda_6\) 的 \(d=2\) 谱层进一步对应于
\[
Z\!\bigl(\mathcal G_6^{\mathrm{bin}}\bigr)
\cong
(\ZZ/2\ZZ)^3_{\partial}\oplus(\ZZ/2\ZZ)^5_{\mathrm{int}}
\subset
\bigl(\mathcal G_6^{\mathrm{bin}}\bigr)^{\mathrm{ab}}
\cong
(\ZZ/2\ZZ)^3\oplus(\ZZ/2\ZZ)^5\oplus(\ZZ/2\ZZ)^{13}.
\]

\paragraph{审计偶数窗上的最小退化扇区、群胚维数、规范群熵与黄金密度分裂}
在定理 \ref{thm:xi-foldbin-minimal-center-valued-index-spectrum-top} 已把纤维多重度场压缩为中心谱变量之后，审计偶数窗
\[
m\in\{6,8,10,12\}
\]
上由最小退化扇区驱动的碰撞、Abel 化、群胚中心与规范群熵还可进一步收束为一组完全对象层结论。沿用定义 \ref{def:fold-bin-degeneracy-sectors}、定理 \ref{thm:fold-bin-min-degeneracy-fib-audited}、推论 \ref{cor:fold-bin-minsector-complement-maxfiber-audited}、定理 \ref{thm:xi-foldbin-collision-probability-normalized-groupoid-dimension}、定理 \ref{thm:xi-foldbin-parity-charge-split-exact-minrank}、命题 \ref{prop:fold-bin-groupoid-algebra-wedderburn}、注记 \ref{rem:fold-bin-gauge-entropy-scale} 与推论 \ref{cor:fold-bin-minsector-density-phi-minus2} 的记号，记
\[
d_*:=d_{\min}(m)=F_{m/2},
\qquad
\mathcal S_{m,*}:=\mathcal S_{m,d_*},
\qquad
r_m:=|X_m\setminus \mathcal S_{m,*}|=D_{2m-2}=F_{m+1}.
\]
为避免与定理 \ref{thm:xi-foldbin-minimal-center-valued-index-spectrum-top} 中的谱顶端记号 \(M_m\) 混淆，记补扇区总质量
\[
T_m:=2^m-F_m d_*,
\qquad
T_m=q_m r_m+s_m,\qquad 0\le s_m<r_m.
\]
并置
\[
A_m:=\CC[\mathcal R_m^{\mathrm{bin}}],
\qquad
A_{m,\min}:=\bigoplus_{w\in\mathcal S_{m,*}}M_{d_*}(\CC),
\qquad
A_{m,\mathrm{out}}:=\bigoplus_{w\in X_m\setminus\mathcal S_{m,*}}
M_{d_m^{\mathrm{bin}}(w)}(\CC).
\]

\begin{theorem}[审计偶数窗连续容量曲线的首折点具有精确 Fibonacci 斜率跳跃]\label{thm:xi-audited-even-first-capacity-kink-fibonacci-jump}
\leanverified{paper\_xi\_audited\_even\_first\_capacity\_kink\_fibonacci\_jump}
沿用定理 \ref{thm:xi-foldbin-continuous-capacity-right-derivative-abelian-center} 的连续容量记号，定义
\[
\mathcal C_m^{\mathrm{cont}}(t):=\sum_{w\in X_m}\min\bigl(d_m^{\mathrm{bin}}(w),t\bigr),
\qquad t\ge 0.
\]
则对每个审计偶数窗 \(m\in\{6,8,10,12\}\)，首折点 \(t=d_*=F_{m/2}\) 附近成立精确分段公式
\[
\boxed{\;
\mathcal C_m^{\mathrm{cont}}(t)=F_{m+2}\,t
\qquad
(0\le t\le d_*),\;}
\]
\[
\boxed{\;
\mathcal C_m^{\mathrm{cont}}(t)=F_m\,d_*+F_{m+1}\,t
\qquad
(d_*\le t\le d_*+1).\;}
\]
从而在首折点处的左右导数满足
\[
\boxed{\;
\partial_t^{-}\mathcal C_m^{\mathrm{cont}}(d_*)=F_{m+2},
\qquad
\partial_t^{+}\mathcal C_m^{\mathrm{cont}}(d_*)=F_{m+1},\;}
\]
并且首个斜率跳跃精确等于
\[
\boxed{\;
\partial_t^{-}\mathcal C_m^{\mathrm{cont}}(d_*)-\partial_t^{+}\mathcal C_m^{\mathrm{cont}}(d_*)
=
F_m.\;}
\]
\end{theorem}
\begin{proof}
由定理 \ref{thm:fold-bin-min-degeneracy-fib-audited}，
\[
d_m^{\mathrm{bin}}(w)\ge d_*
\qquad (w\in X_m),
\]
且
\[
|\mathcal S_{m,*}|=F_m.
\]
因此当 \(0\le t\le d_*\) 时，每一条纤维都尚未饱和，从而
\[
\mathcal C_m^{\mathrm{cont}}(t)
=
\sum_{w\in X_m} t
=
|X_m|\,t
=
F_{m+2}\,t.
\]

再由推论 \ref{cor:fold-bin-minsector-complement-maxfiber-audited}，
\[
|X_m\setminus \mathcal S_{m,*}|=r_m=F_{m+1}.
\]
而对 \(w\in\mathcal S_{m,*}\) 有
\[
d_m^{\mathrm{bin}}(w)=d_*,
\]
对 \(w\in X_m\setminus \mathcal S_{m,*}\) 则由最小退化度定义知
\[
d_m^{\mathrm{bin}}(w)\ge d_*+1.
\]
故当 \(d_*\le t\le d_*+1\) 时，最小退化扇区中的纤维已经饱和，而补扇区中的纤维仍处于线性段，于是
\[
\mathcal C_m^{\mathrm{cont}}(t)
=
\sum_{w\in\mathcal S_{m,*}} d_*
+
\sum_{w\in X_m\setminus \mathcal S_{m,*}} t
=
F_m\,d_*+F_{m+1}\,t.
\]
这就给出了所示两段精确公式。

最后，两段都是线性函数，故首折点处的左右导数分别等于对应斜率
\[
F_{m+2},
\qquad
F_{m+1}.
\]
它们之差为
\[
F_{m+2}-F_{m+1}=F_m,
\]
即首个斜率跳跃恰等于最小退化扇区的 Fibonacci 基数。证毕。
\end{proof}

\begin{theorem}[群胚代数维数与碰撞矩的严格等同及最小扇区平衡下界]\label{thm:xi-audited-even-groupoid-dimension-collision-balanced-lower-bound}
\leanverified{paper\_xi\_audited\_even\_groupoid\_dimension\_collision\_balanced\_lower\_bound}
有严格恒等式
\[
\boxed{\;
\dim_{\CC}A_m
=
\sum_{w\in X_m}\bigl(d_m^{\mathrm{bin}}(w)\bigr)^2
=
4^m\,\mathsf{Col}^{\mathrm{bin}}_m.\;}
\]
进一步，
\[
\boxed{\;
\dim_{\CC}A_m
\ge
F_m d_*^2+(r_m-s_m)q_m^2+s_m(q_m+1)^2.\;}
\]
并且等号成立当且仅当补扇区中的全部纤维大小恰只取 \(q_m\) 与 \(q_m+1\) 两值，其中恰有 \(s_m\) 条取 \(q_m+1\)。
\end{theorem}
\begin{proof}
由定理 \ref{thm:xi-foldbin-collision-probability-normalized-groupoid-dimension}，
\[
\dim_{\CC}A_m
=
\sum_{w\in X_m}\bigl(d_m^{\mathrm{bin}}(w)\bigr)^2
=
4^m\,\mathsf{Col}^{\mathrm{bin}}_m,
\]
第一条结论成立。

再证下界。由定理 \ref{thm:fold-bin-min-degeneracy-fib-audited}，
\[
|\mathcal S_{m,*}|=F_m,
\qquad
d_m^{\mathrm{bin}}(w)=d_*
\quad (w\in\mathcal S_{m,*}),
\]
而由推论 \ref{cor:fold-bin-minsector-complement-maxfiber-audited}，
\[
|X_m\setminus \mathcal S_{m,*}|=r_m.
\]
将补扇区上的纤维大小重记为
\[
e_1,\dots,e_{r_m},
\]
则
\[
e_i\ge d_*+1,
\qquad
\sum_{i=1}^{r_m}e_i=T_m,
\]
并且
\[
\dim_{\CC}A_m
=
F_m d_*^2+\sum_{i=1}^{r_m}e_i^2.
\]
因此只需在固定长度 \(r_m\) 与固定总和 \(T_m\) 的约束下最小化
\(
\sum_i e_i^2
\)。

若某一对 \(a,b\) 满足 \(a\ge b+2\)，则
\[
a^2+b^2-(a-1)^2-(b+1)^2
=
2(a-b)-2
>0.
\]
故在保持总和不变的前提下，用 \((a-1,b+1)\) 代替 \((a,b)\) 会严格降低平方和。不断迭代可知，极小值只能在所有坐标彼此至多相差 \(1\) 时取得。把
\[
T_m=q_m r_m+s_m,\qquad 0\le s_m<r_m
\]
代入，唯一可能的极小配置即为
\[
\underbrace{q_m,\dots,q_m}_{r_m-s_m\text{ 次}},
\underbrace{q_m+1,\dots,q_m+1}_{s_m\text{ 次}}.
\]
于是
\[
\sum_{i=1}^{r_m}e_i^2
\ge
(r_m-s_m)q_m^2+s_m(q_m+1)^2,
\]
从而得到所示下界。

最后说明等号条件。对四个审计偶数窗直接代入可得
\[
q_m\in\{3,5,8,12\},
\qquad
d_*\in\{2,3,5,8\},
\]
故总有 \(q_m\ge d_*+1\)，从而上述平衡配置确为补扇区约束 \(e_i\ge d_*+1\) 下的可行点。于是等号成立当且仅当补扇区中的全部纤维大小恰为 \(q_m\) 与 \(q_m+1\) 两值，且其中有 \(s_m\) 条取 \(q_m+1\)。证毕。
\end{proof}

\begin{theorem}[审计偶数窗上的 Abel 化与群胚代数中心的 Fibonacci 双尺度分裂]\label{thm:xi-audited-even-abel-center-fibonacci-splitting}
\leanverified{paper\_xi\_audited\_even\_abel\_center\_fibonacci\_splitting}
有规范直和分裂
\[
\boxed{\;
\bigl(G_m^{\mathrm{bin}}\bigr)^{\mathrm{ab}}
\cong
(\ZZ/2\ZZ)^{\mathcal S_{m,*}}
\oplus
(\ZZ/2\ZZ)^{X_m\setminus\mathcal S_{m,*}}
\cong
(\ZZ/2\ZZ)^{F_m}\oplus(\ZZ/2\ZZ)^{D_{2m-2}}.\;}
\]
并且
\[
\boxed{\;
Z(A_m)
\cong
Z(A_{m,\min})\oplus Z(A_{m,\mathrm{out}})
\cong
\CC^{F_m}\oplus\CC^{D_{2m-2}}.\;}
\]
其中第一直和因子恰由最小退化扇区 \(\mathcal S_{m,*}\) 提供，第二直和因子恰由其补集提供。
\end{theorem}
\begin{proof}
由定理 \ref{thm:xi-foldbin-parity-charge-split-exact-minrank}，
\[
\bigl(G_m^{\mathrm{bin}}\bigr)^{\mathrm{ab}}
\cong
(\ZZ/2\ZZ)^{N_m},
\qquad
N_m:=\#\{w\in X_m:\ d_m^{\mathrm{bin}}(w)\ge 2\}.
\]
而在审计偶数窗上，定理 \ref{thm:fold-bin-min-degeneracy-fib-audited} 给出
\[
d_{\min}(m)=d_*\ge 2,
\]
故每条纤维都满足 \(d_m^{\mathrm{bin}}(w)\ge 2\)，从而
\[
N_m=|X_m|.
\]
于是
\[
\bigl(G_m^{\mathrm{bin}}\bigr)^{\mathrm{ab}}
\cong
(\ZZ/2\ZZ)^{X_m}.
\]
再由集合分拆
\[
X_m=\mathcal S_{m,*}\sqcup(X_m\setminus\mathcal S_{m,*}),
\]
得到
\[
(\ZZ/2\ZZ)^{X_m}
\cong
(\ZZ/2\ZZ)^{\mathcal S_{m,*}}
\oplus
(\ZZ/2\ZZ)^{X_m\setminus\mathcal S_{m,*}}.
\]
最后由
\[
|\mathcal S_{m,*}|=F_m,
\qquad
|X_m\setminus\mathcal S_{m,*}|=D_{2m-2},
\]
便得第一条分裂。

对中心，由命题 \ref{prop:fold-bin-groupoid-algebra-wedderburn}，
\[
A_m
\cong
\bigoplus_{w\in X_m}M_{d_m^{\mathrm{bin}}(w)}(\CC)
=
A_{m,\min}\oplus A_{m,\mathrm{out}}.
\]
有限维半单代数的中心对直和可加，因此
\[
Z(A_m)\cong Z(A_{m,\min})\oplus Z(A_{m,\mathrm{out}}).
\]
又每个矩阵块的中心均为 \(\CC\)，故
\[
Z(A_{m,\min})\cong \CC^{\mathcal S_{m,*}},
\qquad
Z(A_{m,\mathrm{out}})\cong \CC^{X_m\setminus\mathcal S_{m,*}}.
\]
取维数并代入上述两组基数，即得
\[
Z(A_m)\cong \CC^{F_m}\oplus \CC^{D_{2m-2}}.
\]
证毕。
\end{proof}

\begin{theorem}[规范群熵的最优平衡下界与补扇区熵缺口]\label{thm:xi-audited-even-gauge-entropy-balanced-lower-bound}
\leanverified{paper\_xi\_audited\_even\_gauge\_entropy\_balanced\_lower\_bound}
沿用注记 \ref{rem:fold-bin-gauge-entropy-scale} 的 bits 记号
\[
H_m^{\mathrm{gauge}}
:=
\log_2|G_m^{\mathrm{bin}}|
=
\sum_{w\in X_m}\log_2\bigl(d_m^{\mathrm{bin}}(w)!\bigr).
\]
则有最优下界
\[
\boxed{\;
H_m^{\mathrm{gauge}}
\ge
F_m\log_2(d_*!)
+
(r_m-s_m)\log_2(q_m!)
+
s_m\log_2((q_m+1)!).\;}
\]
并且等号成立当且仅当补扇区中的全部纤维大小恰只取 \(q_m\) 与 \(q_m+1\) 两值。

等价地，若定义补扇区熵缺口
\[
\Delta H_m:=H_m^{\mathrm{gauge}}-F_m\log_2(d_*!),
\]
则
\[
\boxed{\;
\Delta H_m
\ge
(r_m-s_m)\log_2(q_m!)
+
s_m\log_2((q_m+1)!).\;}
\]
\end{theorem}
\begin{proof}
由定义，最小扇区上的每条纤维大小都等于 \(d_*\)，其条数为 \(F_m\)。把补扇区上的纤维大小仍记为
\[
e_1,\dots,e_{r_m},
\qquad
\sum_{i=1}^{r_m}e_i=T_m.
\]
则
\[
H_m^{\mathrm{gauge}}
=
F_m\log_2(d_*!)
+
\sum_{i=1}^{r_m}\log_2(e_i!).
\]

定义
\[
f(n):=\log_2(n!)
\qquad (n\in\ZZ_{\ge 1}).
\]
则
\[
f(n+1)-f(n)=\log_2(n+1),
\]
故差分严格递增，亦即 \(f\) 在正整数上严格离散凸。若某一对 \(a,b\) 满足 \(a\ge b+2\)，则
\[
f(a)+f(b)-f(a-1)-f(b+1)
=
\log_2 a-\log_2(b+1)
>0.
\]
因此在固定总和 \(T_m\) 与固定项数 \(r_m\) 的约束下，\(\sum_i f(e_i)\) 的极小值同样由“尽可能均匀”的配置取得，也就是全部 \(e_i\) 只取
\[
q_m,\qquad q_m+1
\]
两值，其中恰有 \(s_m\) 个取 \(q_m+1\)。由定理 \ref{thm:xi-audited-even-groupoid-dimension-collision-balanced-lower-bound} 证明末尾的可行性说明，该配置在四个审计偶数窗上确为可行。故
\[
\sum_{i=1}^{r_m}\log_2(e_i!)
\ge
(r_m-s_m)\log_2(q_m!)
+
s_m\log_2((q_m+1)!),
\]
代回即得第一条下界。

若取等，则上述离散平衡过程中每一步都不能继续降低函数值，于是全部 \(e_i\) 必彼此至多相差 \(1\)，亦即只可能取 \(q_m\) 与 \(q_m+1\) 两值；反之若补扇区纤维正好取这两值，则显然达到下界。第二条关于 \(\Delta H_m\) 的公式只是把最小扇区项移至左端。证毕。
\end{proof}

\begin{theorem}[子覆盖密度延拓下的 Abel 秩密度与中心密度黄金分裂]\label{thm:xi-audited-even-minsector-algebraic-density-golden-splitting}
\leanverified{paper\_xi\_audited\_even\_minsector\_algebraic\_density\_golden\_splitting}
若“子覆盖同构”恒等式
\[
|\mathcal S_{m,*}|=|X_{m-2}|
\]
沿某个趋于无穷的偶数序列成立，则在同一序列上有
\[
\boxed{\;
\frac{
\operatorname{rank}_{\FF_2}\bigl((\ZZ/2\ZZ)^{\mathcal S_{m,*}}\bigr)
}{
\operatorname{rank}_{\FF_2}\!\bigl((G_m^{\mathrm{bin}})^{\mathrm{ab}}\bigr)
}
=
\frac{\dim_{\CC}Z(A_{m,\min})}{\dim_{\CC}Z(A_m)}
=
\frac{F_m}{F_{m+2}}
\longrightarrow
\varphi^{-2}.\;}
\]
并且补扇区的对应密度满足
\[
\boxed{\;
\frac{
\operatorname{rank}_{\FF_2}\bigl((\ZZ/2\ZZ)^{X_m\setminus\mathcal S_{m,*}}\bigr)
}{
\operatorname{rank}_{\FF_2}\!\bigl((G_m^{\mathrm{bin}})^{\mathrm{ab}}\bigr)
}
=
\frac{\dim_{\CC}Z(A_{m,\mathrm{out}})}{\dim_{\CC}Z(A_m)}
=
\frac{F_{m+1}}{F_{m+2}}
\longrightarrow
\varphi^{-1}.\;}
\]
\end{theorem}
\begin{proof}
由定理 \ref{thm:xi-audited-even-abel-center-fibonacci-splitting}，
\[
\operatorname{rank}_{\FF_2}\bigl((\ZZ/2\ZZ)^{\mathcal S_{m,*}}\bigr)
=
|\mathcal S_{m,*}|,
\qquad
\operatorname{rank}_{\FF_2}\!\bigl((G_m^{\mathrm{bin}})^{\mathrm{ab}}\bigr)
=
|X_m|,
\]
并且
\[
\dim_{\CC}Z(A_{m,\min})=|\mathcal S_{m,*}|,
\qquad
\dim_{\CC}Z(A_m)=|X_m|.
\]
因此最小扇区的 Abel 秩密度与中心密度完全一致，均等于
\[
\frac{|\mathcal S_{m,*}|}{|X_m|}.
\]
在假设下，推论 \ref{cor:fold-bin-minsector-density-phi-minus2} 已证明
\[
\frac{|\mathcal S_{m,*}|}{|X_m|}
\longrightarrow
\varphi^{-2},
\]
故第一条极限成立。

对补扇区，同理有
\[
\frac{
\operatorname{rank}_{\FF_2}\bigl((\ZZ/2\ZZ)^{X_m\setminus\mathcal S_{m,*}}\bigr)
}{
\operatorname{rank}_{\FF_2}\!\bigl((G_m^{\mathrm{bin}})^{\mathrm{ab}}\bigr)
}
=
\frac{|X_m\setminus\mathcal S_{m,*}|}{|X_m|}
=
\frac{F_{m+1}}{F_{m+2}},
\]
以及
\[
\frac{\dim_{\CC}Z(A_{m,\mathrm{out}})}{\dim_{\CC}Z(A_m)}
=
\frac{|X_m\setminus\mathcal S_{m,*}|}{|X_m|}
=
\frac{F_{m+1}}{F_{m+2}}.
\]
再由
\[
\frac{F_{m+1}}{F_{m+2}}
=
1-\frac{F_m}{F_{m+2}}
\longrightarrow
1-\varphi^{-2}
=
\varphi^{-1},
\]
即得第二条极限。证毕。
\end{proof}

\begin{theorem}[审计偶数窗最小退化扇区的模 \texorpdfstring{$6$}{6} 同伦相图]\label{thm:xi-audited-even-minsector-homotopy-mod6-phase}
\leanverified{paper\_xi\_audited\_even\_minsector\_homotopy\_mod6\_phase}
沿用上述记号。对任意 \(x\in \mathcal S_{m,*}\)，记
\[
K_x:=\mathrm{Ind}^{\Delta}(\Gamma_x).
\]
其中 \(\tau(x)\) 沿用定理 \ref{thm:pom-fiber-independence-complex-classification} 的记号。
则在审计偶数窗 \(m\in\{6,8,10,12\}\) 上成立：
\begin{enumerate}
  \item 若 \(m\equiv 0\pmod 6\)，则对所有 \(x\in \mathcal S_{m,*}\) 都有 \(|K_x|\) 可缩；
  \item 若 \(m\equiv 2,4\pmod 6\)，则对所有 \(x\in \mathcal S_{m,*}\) 都有
  \[
  |K_x|\simeq S^{\tau(x)-1}.
  \]
\end{enumerate}
特别地，
\[
m=6,12\ \Longrightarrow\ \mathcal S_{m,*}\ \text{上的纤维独立集复形全部可缩},
\]
\[
m=8,10\ \Longrightarrow\ \mathcal S_{m,*}\ \text{上的纤维独立集复形全部为非可缩球面}.
\]
\end{theorem}
\begin{proof}
对任意 \(x\in \mathcal S_{m,*}\)，在本段的 bin-fold 记号下有
\[
d_m(x)=d_m^{\mathrm{bin}}(x)=d_*=F_{m/2},
\]
其中最后一步由定理 \ref{thm:fold-bin-min-degeneracy-fib-audited} 给出。
另一方面，Fibonacci 奇偶律满足
\[
F_n\equiv 0\pmod 2
\quad\Longleftrightarrow\quad
3\mid n.
\]
因此
\[
d_m(x)\ \text{为偶}
\quad\Longleftrightarrow\quad
3\mid \frac{m}{2}
\quad\Longleftrightarrow\quad
m\equiv 0\pmod 6,
\]
并且
\[
d_m(x)\ \text{为奇}
\quad\Longleftrightarrow\quad
m\equiv 2,4\pmod 6.
\]
再由定理 \ref{thm:pom-fiber-parity-homotopy-equivalence}，
\[
|K_x|\ \text{可缩}
\quad\Longleftrightarrow\quad
d_m(x)\ \text{为偶},
\]
\[
|K_x|\simeq S^{\tau(x)-1}
\quad\Longleftrightarrow\quad
d_m(x)\ \text{为奇},
\]
结论遂得。证毕。
\end{proof}

由此，在审计偶数窗上，最小退化扇区不再只是一个局域计数对象，而是同时控制以下五层结构：
\[
\mathcal S_{m,*}
\Longrightarrow
\mathcal S_{m,*}\ \text{上的}\ |K_x|\ \text{按}\ m\bmod 6\ \text{切换于可缩相与球面相}
\Longrightarrow
\dim_{\CC}A_m=4^m\mathsf{Col}^{\mathrm{bin}}_m
\Longrightarrow
\bigl(G_m^{\mathrm{bin}}\bigr)^{\mathrm{ab}}\ \text{与}\ Z(A_m)\ \text{的 Fibonacci 双裂解}
\Longrightarrow
H_m^{\mathrm{gauge}}\ \text{的最优平衡下界}
\Longrightarrow
(\varphi^{-2},\varphi^{-1})\ \text{型代数密度分裂}.
\]

\input{subsubsec__xi-time-protocol-conclusions-part55ab-minsector-homotopy-phase-transition}

\endinput

\paragraph{审计偶数窗首折点之前的 dyadic 线性段}
连续容量曲线在首个折点处的精确分段公式，进一步把 dyadic 预算下的预言机容量与成功率写成完全显式的前阈值线性段。

\begin{corollary}[审计偶数窗上 dyadic 预算的首段完全线性]\label{cor:xi-audited-even-dyadic-first-capacity-linear-regime}
\leanverified{paper\_xi\_audited\_even\_dyadic\_first\_capacity\_linear\_regime}
对每个审计偶数窗 \(m\in\{6,8,10,12\}\)，定义
\[
C_m(B):=\mathcal C_m^{\mathrm{cont}}(2^B),
\qquad
\mathrm{Succ}_m^{\mathrm{bin}}(B):=2^{-m}C_m(B).
\]
若
\[
0\le B\le \log_2 F_{m/2},
\qquad\text{等价地}\qquad
2^B\le d_{\min}(m)=F_{m/2},
\]
则
\[
C_m(B)=2^B\,F_{m+2},
\qquad
\mathrm{Succ}_m^{\mathrm{bin}}(B)=2^{B-m}F_{m+2}.
\]
因此在首个折点之前，每增加 \(1\) bit，预言机成功率恰好加倍；而一旦越过阈值 \(B=\log_2 F_{m/2}\)，连续容量的边际斜率便由定理 \ref{thm:xi-audited-even-first-capacity-kink-fibonacci-jump} 中的 \(F_{m+2}\) 跳降为 \(F_{m+1}\)。
\end{corollary}
\begin{proof}
由定理 \ref{thm:xi-audited-even-first-capacity-kink-fibonacci-jump}，当
\[
0\le t\le d_{\min}(m)=F_{m/2}
\]
时有
\[
\mathcal C_m^{\mathrm{cont}}(t)=F_{m+2}\,t.
\]
把 \(t=2^B\) 代入，便得到
\[
C_m(B)=\mathcal C_m^{\mathrm{cont}}(2^B)=2^B\,F_{m+2}
\]
对全部 \(0\le B\le \log_2 F_{m/2}\) 成立。两边再除以 \(2^m\)，即得
\[
\mathrm{Succ}_m^{\mathrm{bin}}(B)=2^{B-m}F_{m+2}.
\]
若 \(B+1\le \log_2 F_{m/2}\)，则
\[
\mathrm{Succ}_m^{\mathrm{bin}}(B+1)=2\,\mathrm{Succ}_m^{\mathrm{bin}}(B),
\]
故该段确为精确倍增律。最后，越过首折点后的斜率跳降正是定理 \ref{thm:xi-audited-even-first-capacity-kink-fibonacci-jump} 的右导数结论。证毕。
\end{proof}

\endinput

\paragraph{window-\texorpdfstring{$6$}{6} 微态 Hilbert 空间的可见/隐藏分解与纤维谱几何}
在定理 \ref{thm:xi-foldbin-gauge-representation-distribution-law-m6-spectrum} 的 \(m=6\) 规范群直积分解、定理 \ref{thm:xi-window6-center-exact-splitting-bdry-cyclic} 的 boundary 中心直和因子，以及定理 \ref{thm:xi-window6-cyclic-kernel-bc3-root-datum} 与表 \ref{tab:fold6_b3c3_root_dictionary_b3}、表 \ref{tab:fold6_b3c3_root_dictionary_c3} 的 \(B_3/C_3\) 根字典已经固定之后，还可把 window-\(6\) 的微态层提升为一条严格的有限维表示论--谱几何链：所有根几何只出现于 gauge 固定点商，而隐藏部门完全由各纤维的零和模组成，并在同纤维团图上量子化为三条有限质量通道。

\begin{theorem}[window-\texorpdfstring{$6$}{6} 微态 Hilbert 空间的规范分解]\label{thm:xi-time-part55d-window6-microstate-hilbert-gauge-splitting}
\leanverified{paper\_xi\_time\_part55d\_window6\_microstate\_hilbert\_gauge\_splitting}
沿用定理 \ref{thm:xi-foldbin-gauge-representation-distribution-law-m6-spectrum} 的记号，记
\[
\Omega_6:=\{0,1,\dots,63\},
\qquad
F_w:=\Fold_6^{\mathrm{bin}}{}^{-1}(w),
\qquad
d(w):=|F_w|
\quad (w\in X_6),
\]
并定义
\[
\mathcal H_6:=\ell^2(\Omega_6)
=
\bigoplus_{w\in X_6}\ell^2(F_w).
\]
在
\[
G_6^{\mathrm{bin}}
=
\prod_{w\in X_6}\mathfrak S_{d(w)}
\]
的自然置换表示下，存在规范正交分解
\[
\mathcal H_6=\mathcal H_{\mathrm{vis}}\oplus \mathcal H_{\mathrm{hid}},
\qquad
\mathcal H_{\mathrm{vis}}
:=
\bigoplus_{w\in X_6}\CC \mathbf 1_{F_w}
=
\mathcal H_6^{G_6^{\mathrm{bin}}},
\]
并且
\[
\mathcal H_{\mathrm{hid}}
\cong
\bigoplus_{w:\,d(w)=4}\mathrm{Std}_4^{(w)}
\oplus
\bigoplus_{w:\,d(w)=3}\mathrm{Std}_3^{(w)}
\oplus
\bigoplus_{w:\,d(w)=2}\mathrm{sgn}_2^{(w)}.
\]
因此
\[
\dim \mathcal H_{\mathrm{vis}}=21,
\qquad
\dim \mathcal H_{\mathrm{hid}}=64-21=43,
\]
且
\[
\dim \mathcal H_{\mathrm{hid}}
=
9\cdot 3+4\cdot 2+8\cdot 1
=
27+8+8.
\]
\end{theorem}
\begin{proof}
由于 \(\{F_w:\ w\in X_6\}\) 构成 \(\Omega_6\) 的不交划分，故
\[
\mathcal H_6
=
\bigoplus_{w\in X_6}\ell^2(F_w)
\]
是规范正交直和。对每个 \(w\in X_6\)，都有
\[
\ell^2(F_w)=\CC \mathbf 1_{F_w}\oplus \mathbf 1_{F_w}^{\perp}.
\]
其中 \(\mathfrak S_{d(w)}\) 在 \(\CC \mathbf 1_{F_w}\) 上平凡作用，在 \(\mathbf 1_{F_w}^{\perp}\) 上给出零和超平面表示。于是当 \(d(w)=4\) 时得到三维标准表示 \(\mathrm{Std}_4\)；当 \(d(w)=3\) 时得到二维标准表示 \(\mathrm{Std}_3\)；当 \(d(w)=2\) 时，\(\mathbf 1_{F_w}^{\perp}\) 为一维符号表示 \(\mathrm{sgn}_2\)。

对整个直积群 \(G_6^{\mathrm{bin}}\) 而言，每个因子 \(\mathfrak S_{d(w)}\) 只作用在对应的 \(\ell^2(F_w)\) 上，因此把上述分解逐纤维相加，便得到所示
\[
\mathcal H_6=\mathcal H_{\mathrm{vis}}\oplus \mathcal H_{\mathrm{hid}}
\]
及隐藏部门的不可约分块。

最后，\(\mathcal H_6^{G_6^{\mathrm{bin}}}\) 恰由各纤维上的常数函数组成，故等于 \(\mathcal H_{\mathrm{vis}}\)。维数计算则由推论 \ref{cor:fold-bin-m6-fiber-hist} 的直方图
\[
4^{\times 9},\qquad 3^{\times 4},\qquad 2^{\times 8}
\]
直接给出。证毕。
\end{proof}

\begin{corollary}[window-\texorpdfstring{$6$}{6} 的 \texorpdfstring{$B_3/C_3$}{B3/C3} 根几何只存在于可见固定点商]\label{cor:xi-time-part55d-window6-root-geometry-visible-fixed-quotient}
沿用定理 \ref{thm:xi-time-part55d-window6-microstate-hilbert-gauge-splitting} 的记号，定义
\[
\mathcal H_{\mathrm{vis}}^{\mathrm{cyc}}
:=
\bigoplus_{w\in X_6^{\mathrm{cyc}}}\CC \mathbf 1_{F_w},
\qquad
\mathcal H_{\mathrm{vis}}^{\mathrm{bdry}}
:=
\bigoplus_{w\in X_6^{\mathrm{bdry}}}\CC \mathbf 1_{F_w}.
\]
则
\[
\mathcal H_{\mathrm{vis}}
=
\mathcal H_{\mathrm{vis}}^{\mathrm{cyc}}
\oplus
\mathcal H_{\mathrm{vis}}^{\mathrm{bdry}}
\cong
\CC^{18}\oplus \CC^3.
\]
并且表 \ref{tab:fold6_b3c3_root_dictionary_b3} 与表 \ref{tab:fold6_b3c3_root_dictionary_c3} 的 \(B_3/C_3\) 字典仅定义在 \(\mathcal H_{\mathrm{vis}}^{\mathrm{cyc}}\) 的自然基上。换言之，window-\(6\) 的根系几何并非微态层上的结构，而是 gauge 固定点商中的可见几何。
\end{corollary}
\begin{proof}
由定理 \ref{thm:xi-window6-cyclic-kernel-bc3-root-datum}，
\[
X_6=X_6^{\mathrm{cyc}}\sqcup X_6^{\mathrm{bdry}},
\qquad
|X_6^{\mathrm{cyc}}|=18,
\qquad
|X_6^{\mathrm{bdry}}|=3.
\]
因此 \(\mathcal H_{\mathrm{vis}}\) 按 cyclic/boundary 精确裂解为
\[
\mathcal H_{\mathrm{vis}}^{\mathrm{cyc}}\oplus \mathcal H_{\mathrm{vis}}^{\mathrm{bdry}}
\cong
\CC^{18}\oplus \CC^3.
\]
另一方面，表 \ref{tab:fold6_b3c3_root_dictionary_b3} 与表 \ref{tab:fold6_b3c3_root_dictionary_c3} 只对 \(X_6^{\mathrm{cyc}}\) 的 \(18\) 个稳定词给出确定性根字典，而不涉及 boundary 三词。故根几何完全落在 \(\mathcal H_{\mathrm{vis}}^{\mathrm{cyc}}\) 上。证毕。
\end{proof}

\begin{theorem}[规范作用像代数与交换子代数的双侧刚性]\label{thm:xi-time-part55d-window6-action-image-commutant-rigidity}
设
\[
\rho:\CC[G_6^{\mathrm{bin}}]\longrightarrow \operatorname{End}(\mathcal H_6)
\]
为自然置换表示的线性延拓。则其像代数满足
\[
\rho\!\bigl(\CC[G_6^{\mathrm{bin}}]\bigr)
\cong
\CC
\oplus
M_3(\CC)^{\oplus 9}
\oplus
M_2(\CC)^{\oplus 4}
\oplus
\CC^{\oplus 8},
\]
而其交换子代数满足
\[
\rho\!\bigl(\CC[G_6^{\mathrm{bin}}]\bigr)'
\cong
M_{21}(\CC)
\oplus
\CC^{\oplus 9}
\oplus
\CC^{\oplus 4}
\oplus
\CC^{\oplus 8}
\cong
M_{21}(\CC)\oplus \CC^{\oplus 21}.
\]
特别地，\(T\in \operatorname{End}(\mathcal H_6)\) 为自伴且与 \(G_6^{\mathrm{bin}}\) 作用可交换，当且仅当
\[
T
=
A_{\mathrm{vis}}
\oplus
\bigoplus_{w\in X_6}\lambda_w I_{\mathbf 1_{F_w}^{\perp}},
\qquad
A_{\mathrm{vis}}=A_{\mathrm{vis}}^*\in \operatorname{End}(\mathcal H_{\mathrm{vis}}),
\quad
\lambda_w\in \RR.
\]
因此，一切与 gauge 作用等变的非平凡耦合都被压缩在 \(21\) 维可见商 \(\mathcal H_{\mathrm{vis}}\) 中；隐藏部门上既不存在跨纤维耦合，也不存在可见--隐藏混合耦合。
\end{theorem}
\begin{proof}
由定理 \ref{thm:xi-time-part55d-window6-microstate-hilbert-gauge-splitting}，
\[
\mathcal H_6
\cong
\mathbf 1^{\oplus 21}
\oplus
\bigoplus_{w:\,d(w)=4}\mathrm{Std}_4^{(w)}
\oplus
\bigoplus_{w:\,d(w)=3}\mathrm{Std}_3^{(w)}
\oplus
\bigoplus_{w:\,d(w)=2}\mathrm{sgn}_2^{(w)}.
\]
对每个 \(w\)，隐藏块 \(\mathbf 1_{F_w}^{\perp}\) 上只有第 \(w\) 个因子 \(\mathfrak S_{d(w)}\) 非平凡作用，其余因子都平凡。因此每个隐藏块都是 \(G_6^{\mathrm{bin}}\) 的不可约表示；若 \(w\neq w'\)，则把表示限制到第 \(w\) 个因子即可看出：\(\mathbf 1_{F_w}^{\perp}\) 上该因子非平凡，而 \(\mathbf 1_{F_{w'}}^{\perp}\) 上该因子平凡，故这些隐藏块两两不同构。

另一方面，Burnside 定理给出
\[
\CC[\mathfrak S_4]\big|_{\mathrm{Std}_4}\cong M_3(\CC),
\qquad
\CC[\mathfrak S_3]\big|_{\mathrm{Std}_3}\cong M_2(\CC),
\qquad
\CC[\mathfrak S_2]\big|_{\mathrm{sgn}_2}\cong \CC.
\]
而 \(21\) 条常数线共同构成平凡表示的 \(21\) 重等型分量，群代数在其上只通过标量作用。于是像代数正是
\[
\CC
\oplus
M_3(\CC)^{\oplus 9}
\oplus
M_2(\CC)^{\oplus 4}
\oplus
\CC^{\oplus 8}.
\]

再由 Schur 引理，交换子在平凡表示的 \(21\) 重等型块上给出全部矩阵代数
\[
\operatorname{End}(\mathcal H_{\mathrm{vis}})\cong M_{21}(\CC),
\]
而在每个 multiplicity-one 的隐藏不可约块上只能给出标量。故
\[
\rho\!\bigl(\CC[G_6^{\mathrm{bin}}]\bigr)'
\cong
M_{21}(\CC)\oplus \CC^{\oplus 21}.
\]
最后，自伴条件只要求 \(A_{\mathrm{vis}}\) 自伴且各标量 \(\lambda_w\) 取实数。由不同不可约分量之间不存在非零交织算子，隐藏部门上不可能出现跨纤维项，也不可能与可见部门混合。证毕。
\end{proof}

\begin{theorem}[粗粒化算子的三值奇异谱]\label{thm:xi-time-part55d-window6-coarse-graining-three-singular-values}
\leanverified{paper\_xi\_time\_part55d\_window6\_coarse\_graining\_three\_singular\_values}
在 \(\ell^2(X_6)\) 的标准正交基 \((e_w)_{w\in X_6}\) 下，定义粗粒化算子
\[
B:\mathcal H_6\longrightarrow \ell^2(X_6),
\qquad
(Bf)(w):=\sum_{n\in F_w}f(n),
\]
并记
\[
D:=\operatorname{diag}(d(w))_{w\in X_6}.
\]
则成立：
\begin{enumerate}
  \item
  \[
  B^*B=\bigoplus_{w\in X_6}J_{d(w)},
  \]
  其中 \(J_d\) 为 \(d\times d\) 全 \(1\) 矩阵；
  \item \(B\) 的非零奇异值恰为
  \[
  \{2^{(9)},\ \sqrt3^{(4)},\ \sqrt2^{(8)}\};
  \]
  \item
  \[
  \ker B=\mathcal H_{\mathrm{hid}},
  \qquad
  \operatorname{rank}B=21;
  \]
  \item 若定义
  \[
  \widetilde B:=D^{-1/2}B,
  \]
  则
  \[
  \widetilde B\widetilde B^*=I_{\ell^2(X_6)},
  \qquad
  \widetilde B^*\widetilde B=P_{\mathrm{vis}},
  \]
  其中 \(P_{\mathrm{vis}}\) 为到 \(\mathcal H_{\mathrm{vis}}\) 的正交投影。特别地，\(\ell^2(X_6)\) 与 \(\mathcal H_{\mathrm{vis}}\) 规范等距同构。
\end{enumerate}
\end{theorem}
\begin{proof}
对每个 \(w\in X_6\)，算子 \(B\) 在 \(\ell^2(F_w)\) 上的限制只是求和泛函：
\[
B|_{\ell^2(F_w)}(f)
=
\left(\sum_{n\in F_w}f(n)\right)e_w.
\]
故
\[
B^*e_w=\mathbf 1_{F_w},
\]
从而 \(B^*B\) 在第 \(w\) 个纤维块上正是 \(J_{d(w)}\)，这就得到第 \(1\) 条。

而 \(J_d\) 的特征值为 \(d\)（重数 \(1\)）与 \(0\)（重数 \(d-1\)），所以 \(B\) 的非零奇异值恰为 \(\sqrt{d(w)}\)。由推论 \ref{cor:fold-bin-m6-fiber-hist} 的直方图，\(d(w)=4,3,2\) 的纤维个数分别为 \(9,4,8\)，于是得到第 \(2\) 条。

同样，\(J_d\) 的零空间正是零和超平面 \(\mathbf 1_{F_w}^{\perp}\)，故
\[
\ker B
=
\bigoplus_{w\in X_6}\mathbf 1_{F_w}^{\perp}
=
\mathcal H_{\mathrm{hid}}.
\]
由
\[
\dim \mathcal H_6=64,
\qquad
\dim \ker B=43
\]
立得 \(\operatorname{rank}B=21\)，从而第 \(3\) 条成立。

最后，直接计算得
\[
BB^*=D,
\]
故
\[
\widetilde B\widetilde B^*
=
D^{-1/2}BB^*D^{-1/2}
=
I_{\ell^2(X_6)}.
\]
另一方面，\(\widetilde B^*\widetilde B\) 在第 \(w\) 个纤维块上等于
\[
d(w)^{-1}J_{d(w)},
\]
这正是到 \(\CC \mathbf 1_{F_w}\) 的正交投影。逐块相加便得到
\[
\widetilde B^*\widetilde B=P_{\mathrm{vis}}.
\]
因此 \(\widetilde B\) 把 \(\mathcal H_{\mathrm{vis}}\) 等距识别为 \(\ell^2(X_6)\)。证毕。
\end{proof}

\begin{theorem}[同纤维团图的质量量子化与热核三指数分解]\label{thm:xi-time-part55d-window6-fiber-clique-mass-quantization}
在 \(\Omega_6\) 上定义图 \(\Gamma_6\)：若且唯若 \(u\neq v\) 且
\[
\Fold_6^{\mathrm{bin}}(u)=\Fold_6^{\mathrm{bin}}(v),
\]
则在 \(u,v\) 之间连一条边。记 \(A_{\Gamma_6}\) 为其邻接矩阵，\(L_{\Gamma_6}\) 为其图拉普拉斯。则
\[
\Gamma_6\cong 9K_4\sqcup 4K_3\sqcup 8K_2.
\]
从而
\[
\operatorname{Spec}(A_{\Gamma_6})
=
\{3^{(9)},\ 2^{(4)},\ 1^{(8)},\ (-1)^{(43)}\},
\]
\[
\operatorname{Spec}(L_{\Gamma_6})
=
\{0^{(21)},\ 4^{(27)},\ 3^{(8)},\ 2^{(8)}\}.
\]
若 \(P_0,P_2,P_3,P_4\) 为 \(L_{\Gamma_6}\) 对应于特征值 \(0,2,3,4\) 的谱投影，则
\[
P_0=P_{\mathrm{vis}},
\qquad
\mathcal H_{\mathrm{hid}}
=
P_2\mathcal H_6\oplus P_3\mathcal H_6\oplus P_4\mathcal H_6,
\]
并且
\[
\dim P_2\mathcal H_6=8,
\qquad
\dim P_3\mathcal H_6=8,
\qquad
\dim P_4\mathcal H_6=27.
\]
进一步，对任意 \(t\ge 0\) 都有精确热核分解
\[
e^{-tL_{\Gamma_6}}
=
P_{\mathrm{vis}}
+
e^{-2t}P_2
+
e^{-3t}P_3
+
e^{-4t}P_4.
\]
因此隐藏部门不含零模，而严格量子化为三条有限质量通道 \(2,3,4\)。
\end{theorem}
\begin{proof}
按定义，\(\Gamma_6\) 的每个连通分支恰是某条纤维 \(F_w\) 上的完全图 \(K_{d(w)}\)。由推论 \ref{cor:fold-bin-m6-fiber-hist}，window-\(6\) 的纤维直方图为
\[
4^{\times 9},\qquad 3^{\times 4},\qquad 2^{\times 8},
\]
故
\[
\Gamma_6\cong 9K_4\sqcup 4K_3\sqcup 8K_2.
\]

对完全图 \(K_d\)，邻接谱为
\[
\{(d-1)^{(1)},\ (-1)^{(d-1)}\},
\]
图拉普拉斯谱为
\[
\{0^{(1)},\ d^{(d-1)}\}.
\]
逐分支相加即得所示总谱。

特别地，\(L_{\Gamma_6}\) 的零特征空间正是每个连通分支上的常数函数空间，因此恰为
\[
\bigoplus_{w\in X_6}\CC \mathbf 1_{F_w}
=
\mathcal H_{\mathrm{vis}}.
\]
相应的正特征空间则是各分支上的零和空间，因而等于 \(\mathcal H_{\mathrm{hid}}\)。其中 \(K_2\) 贡献 \(8\) 维特征值 \(2\) 空间，\(K_3\) 贡献 \(4\cdot 2=8\) 维特征值 \(3\) 空间，\(K_4\) 贡献 \(9\cdot 3=27\) 维特征值 \(4\) 空间。热核公式最后由谱定理直接给出。证毕。
\end{proof}

\begin{conjecture}[六顶点/GFF 极限中的可见--隐藏严格分离]\label{conj:xi-time-part55d-window6-sixvertex-gff-visible-hidden-separation}
沿用定理 \ref{thm:xi-time-part55d-window6-fiber-clique-mass-quantization} 的谱投影 \(P_{\mathrm{vis}},P_2,P_3,P_4\)。设某一临界六顶点尺度极限下，高度函数在适当归一化后收敛到标量高斯自由场。则任意随网格尺度 \(\delta\downarrow 0\) 变化的 window-\(6\) 局部观测量提升
\[
\mathcal O_\delta\in \mathcal H_6
\]
都应满足正交分解
\[
\mathcal O_\delta
=
P_{\mathrm{vis}}\mathcal O_\delta
+
P_2\mathcal O_\delta
+
P_3\mathcal O_\delta
+
P_4\mathcal O_\delta,
\]
并且：
\begin{enumerate}
  \item 只有 \(P_{\mathrm{vis}}\mathcal O_\delta\) 可能产生临界长程极限；
  \item \(P_2\mathcal O_\delta\)、\(P_3\mathcal O_\delta\)、\(P_4\mathcal O_\delta\) 全部只贡献局部且紧的有限质量修正；
  \item window-\(6\) 的隐藏部门不可能再生成第二个独立的无质量 Gaussian 场；
  \item 表 \ref{tab:fold6_b3c3_root_dictionary_b3} 与表 \ref{tab:fold6_b3c3_root_dictionary_c3} 所编码的 \(B_3/C_3\) 根系几何，仅是可见固定点商上的组合记账，而不是新的连续自由场来源。
\end{enumerate}
\end{conjecture}

\begin{remark}
该猜想的要点不在于再次排除额外高斯模，而在于把全部连续极限障碍精确压缩为隐藏部门的三条离散质量 \(2,3,4\)。若其成立，则 window-\(6\) 与六顶点/GFF 的接口将被压缩为一个单一命题：临界性只驻留于 \(21\) 维可见商，而全部隐藏模式只能以有限质量方式附着其上。
\end{remark}

\endinput

\paragraph{同余类残差、审计容量硬界、可见算术计数与零温地基层}
前述 Dirichlet 扭曲谱、审计偶数窗上的 dyadic 容量、有限分辨率可见算术以及输出位势的零温极限，还可继续压缩为下列五条直接后果。周期层的同余类平方根障碍已由定理 \ref{thm:xi-dirichlet-mode-forces-congruence-residual} 给出；下条表明，经 Witt/M\"obius 投影之后，这一障碍在 primitive 层并不会消失；最后一条则把地基层熵、端点极点尺度与零温矩阵的谱半径压缩为同一条平方律。

\begin{theorem}[Witt/M\"obius 投影下的 primitive 同余类平方根障碍]\label{thm:xi-primitive-congruence-residual-nonrh}
沿用命题 \ref{prop:real-input-40-output-dirichlet-nonrh} 与定理 \ref{thm:xi-dirichlet-mode-forces-congruence-residual} 的记号。对每个模数 \(m\ge 2\) 与剩余类 \(a\in\ZZ/m\ZZ\)，定义 primitive 层的同余计数
\[
P_n(a;m):=\#\{\mathcal O:\ \mathcal O\text{ 为长度 }n\text{ 的 primitive 轨道},\ E(\mathcal O)\equiv a\pmod m\},
\qquad
P_n:=\sum_{a\in\ZZ/m\ZZ}P_n(a;m).
\]
若
\[
\theta_m:=\max_{1\le j\le m-1}\frac{\log\rho_{m,j}}{\log\lambda}>\frac12,
\qquad
\lambda=\rho(M(1))=\varphi^2,
\]
则存在某个剩余类 \(a_\ast\in\ZZ/m\ZZ\) 使
\[
\limsup_{n\to\infty}
\frac{1}{n\log\lambda}
\log\Bigl|P_n(a_\ast;m)-\frac{P_n}{m}\Bigr|
\ge
\theta_m
>
\frac12.
\]
因此，在该模数下，primitive 层的同余类分布同样不可能满足任何 RH 形的平方根指数误差界。
\end{theorem}
\begin{proof}
取 \(j_\ast\in\{1,\dots,m-1\}\) 使
\[
\rho_{m,j_\ast}=\lambda^{\theta_m}=:\rho>\lambda^{1/2}.
\]
定义 primitive Fourier 系数
\[
\widehat P_n(j):=\sum_{a\in\ZZ/m\ZZ}P_n(a;m)\omega_m^{ja}.
\]
由命题 \ref{prop:real-input-40-output-dirichlet-nonrh} 中的 Witt/M\"obius 反演，
\[
\widehat P_n(j_\ast)
=
p_n(\omega_m^{j_\ast})
=
\frac{1}{n}\sum_{d\mid n}\mu(d)\,a_{n/d}(\omega_m^{j_\ast d}).
\]
另一方面，定理 \ref{thm:xi-dirichlet-mode-forces-congruence-residual} 的证明已给出
\[
\limsup_{n\to\infty}|a_n(\omega_m^{j_\ast})|^{1/n}=\rho,
\]
故存在子序列 \(n_r\to\infty\) 满足
\[
|a_{n_r}(\omega_m^{j_\ast})|=\rho^{\,n_r+o(n_r)}.
\]

现估计 Möbius 反演中的其余约数项。对任意 \(d\ge 2\)，由逐项取绝对值与 Perron 单调性，
\[
\rho\!\bigl(M(\omega_m^{j_\ast d})\bigr)\le \rho(M(1))=\lambda,
\]
从而
\[
a_{n_r/d}(\omega_m^{j_\ast d})=O(\lambda^{n_r/d})=O(\lambda^{n_r/2})=o(\rho^{n_r}),
\]
因为 \(\rho>\lambda^{1/2}\)。又约数函数满足 \(\tau(n_r)=e^{o(n_r)}\)，故
\[
\sum_{\substack{d\mid n_r\\ d\ge 2}}\mu(d)\,a_{n_r/d}(\omega_m^{j_\ast d})=o(\rho^{n_r}).
\]
于是沿该子序列有
\[
\widehat P_{n_r}(j_\ast)
=
\frac{1}{n_r}a_{n_r}(\omega_m^{j_\ast})+o\!\left(\frac{\rho^{n_r}}{n_r}\right),
\]
进而
\[
\limsup_{n\to\infty}|\widehat P_n(j_\ast)|^{1/n}\ge \rho.
\]

记 primitive 层的同余残差为
\[
b_n^{\mathrm{prim}}(a;m):=P_n(a;m)-\frac{P_n}{m}.
\]
则对 \(j\neq 0\) 有
\[
\sum_{a\in\ZZ/m\ZZ} b_n^{\mathrm{prim}}(a;m)\,\omega_m^{ja}
=
\widehat P_n(j),
\]
因为 \(\sum_a\omega_m^{ja}=0\)。故
\[
|\widehat P_n(j_\ast)|
\le
\sum_{a\in\ZZ/m\ZZ}|b_n^{\mathrm{prim}}(a;m)|
\le
m\max_{a\in\ZZ/m\ZZ}|b_n^{\mathrm{prim}}(a;m)|.
\]
取 \(n\to\infty\) 的上极限并开 \(n\) 次方，得到
\[
\limsup_{n\to\infty}
\left(
\max_{a\in\ZZ/m\ZZ}\Bigl|P_n(a;m)-\frac{P_n}{m}\Bigr|
\right)^{1/n}
\ge
\rho=\lambda^{\theta_m}.
\]
由于剩余类只有有限个，必存在某个 \(a_\ast\in\ZZ/m\ZZ\) 使相应单类残差的上极限至少为 \(\lambda^{\theta_m}\)。两边取对数并除以 \(n\log\lambda\)，即得所示结论。证毕。
\end{proof}

\leanverified{paper\_xi\_foldbin\_audited\_window\_oracle\_hard\_cap}
\begin{theorem}[审计偶数窗上最小退化扇区强制预言机成功率硬上界]\label{thm:xi-foldbin-audited-window-oracle-hard-cap}
沿用定理 \ref{thm:xi-foldbin-oracle-capacity-central-spectral-truncation} 的记号，并定义
\[
\mathrm{Succ}_m^{\mathrm{bin}}(B):=\frac{C_m(B)}{2^m},
\qquad
(x)_+:=\max\{x,0\}.
\]
若 \(m\in\{6,8,10,12\}\)，则对任意整数 \(B\ge 0\) 都有
\[
\boxed{\;
\mathrm{Succ}_m^{\mathrm{bin}}(B)
\le
1-\frac{F_m}{2^m}\bigl(F_{m/2}-2^B\bigr)_+.\;}
\]
特别地，当 \(B<\log_2 F_{m/2}\) 时，
\[
1-\mathrm{Succ}_m^{\mathrm{bin}}(B)
\ge
\frac{F_m}{2^m}\bigl(F_{m/2}-2^B\bigr)
>
0.
\]
\end{theorem}
\begin{proof}
由定理 \ref{thm:fold-bin-min-degeneracy-fib-audited}，当 \(m\in\{6,8,10,12\}\) 时，恰有 \(F_m\) 条纤维达到最小大小
\[
d_{\min}(m)=F_{m/2}.
\]
在这 \(F_m\) 条纤维上，容量公式中的每一项至多贡献 \(2^B\) 个可恢复微态；因此它们总共至多贡献
\[
F_m\min(F_{m/2},2^B).
\]
其余纤维即使全部完全恢复，也至多贡献
\[
2^m-F_mF_{m/2}
\]
个微态。于是
\[
\begin{aligned}
C_m(B)
&\le
\bigl(2^m-F_mF_{m/2}\bigr)+F_m\min(F_{m/2},2^B)\\
&=
2^m-F_m\bigl(F_{m/2}-2^B\bigr)_+.
\end{aligned}
\]
两边除以 \(2^m\) 即得所示成功率上界。证毕。
\end{proof}

\begin{theorem}[有限分辨率可见算术的单位元与幂等元计数闭式]\label{thm:xi-visible-arithmetic-unit-idempotent-count}
\leanverified{paper\_xi\_visible\_arithmetic\_unit\_idempotent\_count}
沿用定理 \ref{thm:fold-congruence-mod-semirings} 的有限分辨率可见算术记号，记
\[
R_m^{\mathrm{vis}}:=\Omega_m/\!\equiv_m
\cong
\ZZ/(F_{m+2}\ZZ).
\]
则其乘法结构满足
\[
\boxed{\;
\#\bigl(R_m^{\mathrm{vis}}\bigr)^\times
=
\phi_{\mathrm{Euler}}(F_{m+2}),\;}
\]
\[
\boxed{\;
\#\{e\in R_m^{\mathrm{vis}}:\ e^2=e\}
=
2^{\omega(F_{m+2})},\;}
\]
其中 \(\phi_{\mathrm{Euler}}\) 为 Euler totient，\(\omega(N)\) 表示 \(N\) 的不同素因子个数。
\end{theorem}
\begin{proof}
由定理 \ref{thm:fold-congruence-mod-semirings}，
\[
R_m^{\mathrm{vis}}\cong \ZZ/(F_{m+2}\ZZ)
\]
作为半环同构成立。故其乘法单位元类与模 \(F_{m+2}\) 的单位元一一对应，数量正是
\[
\phi_{\mathrm{Euler}}(F_{m+2}).
\]

再把 Fibonacci 模数分解为
\[
F_{m+2}=\prod_{i=1}^{r}p_i^{\nu_i},
\qquad
r=\omega(F_{m+2}).
\]
中国剩余定理给出
\[
\ZZ/(F_{m+2}\ZZ)\cong \prod_{i=1}^{r}\ZZ/(p_i^{\nu_i}\ZZ).
\]
在每个局部因子 \(\ZZ/(p_i^{\nu_i}\ZZ)\) 中，幂等方程 \(e^2=e\) 只有 \(e=0,1\) 两个解；因此全局幂等元可在每个坐标独立选取 \(0\) 或 \(1\)，总数恰为 \(2^r\)。这与定理 \ref{thm:xi-fold-congruence-semiring-singularity-nilpotent-profile} 中给出的幂等元计数完全一致。证毕。
\end{proof}

\begin{theorem}[Fibonacci 同余拓扑与 profinite 拓扑完全一致]\label{thm:xi-visible-arithmetic-fibonacci-adic-profinite-coincidence}
\leanverified{paper\_xi\_visible\_arithmetic\_fibonacci\_adic\_profinite\_coincidence}
沿用定理 \ref{thm:fold-congruence-mod-semirings} 的记号。令
\[
\mathcal N_F:=\{F_{m+2}\ZZ:\ m\ge 1\}
\]
并令 \(\tau_F\) 为使 \(\mathcal N_F\) 中每个理想都成为 \(0\) 邻域的最弱线性拓扑。则 \(\tau_F\) 与 \(\ZZ\) 上的 profinite 拓扑 \(\tau_{\mathrm{pro}}\) 完全一致。因而 \(\tau_F\)-完备化满足规范拓扑环同构
\[
\boxed{\;
\widehat{\ZZ}_F\cong \widehat{\ZZ}.\;}
\]
\end{theorem}
\begin{proof}
先证 \(\tau_F\preceq\tau_{\mathrm{pro}}\)。对每个 \(m\ge 1\)，理想
\[
F_{m+2}\ZZ
\]
本身就是 profinite 邻域基 \(\{N\ZZ:\ N\ge 1\}\) 中的一项，故由 \(\mathcal N_F\) 生成的拓扑不强于 profinite 拓扑。

再证反向包含。任取 \(N\ge 1\)。考虑 Fibonacci 递推在有限集
\[
(\ZZ/N\ZZ)^2
\]
上诱导的状态演化
\[
(F_k,F_{k+1})\longmapsto (F_{k+1},F_{k+2}).
\]
由于状态空间有限，轨道必最终周期；而初态 \((F_0,F_1)=(0,1)\) 在周期中必再现。故存在某个 \(L=L(N)\ge 1\) 使
\[
(F_L,F_{L+1})\equiv (0,1)\pmod N.
\]
于是 \(N\mid F_L\)。若 \(N>1\)，则必有 \(L\ge 3\)，从而可写 \(L=m+2\)；对 \(N=1\) 则结论平凡。因而对任意 \(N\ge 1\) 都存在 \(m\ge 1\) 使
\[
F_{m+2}\ZZ\subseteq N\ZZ.
\]
这说明每个 profinite 邻域都包含某个 Fibonacci 邻域，故 \(\tau_{\mathrm{pro}}\preceq\tau_F\)。

综上 \(\tau_F=\tau_{\mathrm{pro}}\)。两种线性拓扑既然一致，其 Hausdorff 完备化作为同一连续有限商系统的逆极限必自然同构，即
\[
\widehat{\ZZ}_F\cong \widehat{\ZZ}.
\]
证毕。
\end{proof}

\begin{corollary}[可见同余半环族对全部有限模算术的共终性]\label{cor:xi-visible-arithmetic-fibonacci-cofinal-quotients}
\leanverified{paper\_xi\_visible\_arithmetic\_fibonacci\_cofinal\_quotients}
沿用定理 \ref{thm:xi-visible-arithmetic-fibonacci-adic-profinite-coincidence} 与定理 \ref{thm:fold-congruence-mod-semirings} 的记号。对任意正整数 \(N\)，存在某个 \(m\ge 1\) 与满射半环同态
\[
\Psi_{m,N}:R_m^{\mathrm{vis}}=\Omega_m/\!\equiv_m\longtwoheadrightarrow \ZZ/(N\ZZ).
\]
等价地，有限半环族
\[
\{\,\Omega_m/\!\equiv_m\,\}_{m\ge 1}
\]
在所有整数商环 \(\ZZ/(N)\) 的意义下是共终的：每个 \(\ZZ/(N)\) 都是某个 \(\Omega_m/\!\equiv_m\) 的商。
\end{corollary}
\begin{proof}
由定理 \ref{thm:xi-visible-arithmetic-fibonacci-adic-profinite-coincidence} 的证明，对任意 \(N\ge 1\) 可取 \(m\ge 1\) 使
\[
N\mid F_{m+2}.
\]
于是标准降模映射给出满射环同态
\[
\rho_{m,N}:\ZZ/(F_{m+2}\ZZ)\longtwoheadrightarrow \ZZ/(N\ZZ).
\]
再由定理 \ref{thm:fold-congruence-mod-semirings}，
\[
R_m^{\mathrm{vis}}=\Omega_m/\!\equiv_m\cong \ZZ/(F_{m+2}\ZZ),
\]
将 \(\rho_{m,N}\) 经此同构拉回，即得所需的满射半环同态 \(\Psi_{m,N}\)。证毕。
\end{proof}

\begin{theorem}[输出位势零温极限的正熵四态地基层与饱和密度上界]\label{thm:xi-real-input-40-output-freezing-positive-entropy-halffill}
\leanverified{paper\_xi\_real\_input\_40\_output\_freezing\_positive\_entropy\_halffill}
沿用命题 \ref{prop:real-input-40-pressure-freezing}、推论 \ref{cor:real-input-40-ground-entropy} 与定理 \ref{thm:killo-real-input-40-positive-entropy-freezing} 的记号。记
\[
P(\theta):=\log\lambda(e^\theta).
\]
则端点斜率满足
\[
\lim_{\theta\to-\infty}P'(\theta)=0,
\qquad
\lim_{\theta\to+\infty}P'(\theta)=\frac12.
\]
并且零温缩放极限对应于四态一步 SFT \(\Sigma_\infty\)，其转移矩阵为
\[
B_\infty=
\begin{pmatrix}
0&1&1&0\\
0&0&1&0\\
0&0&3&1\\
1&0&0&3
\end{pmatrix},
\]
其地基层熵满足
\[
h_{\mathrm{ground}}
=
\log c
=
\frac12\log\rho(B_\infty)
>
0.
\]
因此，极端正偏置下的冻结端面不可能塌缩为有限个周期轨道的并；同时，可见输出 \(1\) 的典型密度始终被压在区间 \([0,\tfrac12]\) 内。
\end{theorem}
\begin{proof}
命题 \ref{prop:real-input-40-pressure-freezing} 已给出
\[
\lim_{\theta\to-\infty}P'(\theta)=0,
\qquad
\lim_{\theta\to+\infty}P'(\theta)=\frac12,
\]
故极端偏置下的典型输出 \(1\) 密度只能落在 \([0,\tfrac12]\) 内。

另一方面，推论 \ref{cor:real-input-40-ground-entropy} 与定理 \ref{thm:killo-real-input-40-positive-entropy-freezing} 说明：饱和端点对应的零温地基并非有限个最大平均能量周期轨道，而是由上式 \(B_\infty\) 所定义的四态 SFT 支配，且其地基层熵满足
\[
h_{\mathrm{ground}}=\log c=\frac12\log\rho(B_\infty)>0.
\]
若冻结端面仅由有限个周期轨道支撑，则其拓扑熵必为 \(0\)，与 \(h_{\mathrm{ground}}>0\) 矛盾。故极端正偏置极限不可能塌缩为有限周期轨道并。证毕。
\end{proof}

\begin{corollary}[正负偏置零温极限分别冻结于 half-filling 与空相]\label{cor:xi-real-input-40-output-freezing-half-filling-empty-phase}
沿用定理 \ref{thm:xi-real-input-40-output-freezing-positive-entropy-halffill} 的记号。记 \(\mu_\theta\) 为输出位势 \(P(\theta)\) 所对应的平衡态族，并记
\[
\mu_\infty^{+}:=\lim_{\theta\to+\infty}\mu_\theta,
\qquad
\mu_\infty^{-}:=\lim_{\theta\to-\infty}\mu_\theta
\]
为两端零温极限测度。则对观测量 \(\mathbf 1_{\mathrm{out}=1}\) 有
\[
\boxed{\;
\int \mathbf 1_{\mathrm{out}=1}\,\dd\mu_\infty^{+}
=
\frac12,\qquad
\int \mathbf 1_{\mathrm{out}=1}\,\dd\mu_\infty^{-}
=
0.\;}
\]
换言之，正偏置零温极限并不饱和到全占据，而是冻结于严格的 half-filling 地基态；负偏置零温极限则冻结于空相。
\end{corollary}
\begin{proof}
对一参数平衡态族，压力导数给出相应观测量的平衡态期望，即
\[
P'(\theta)=\int \mathbf 1_{\mathrm{out}=1}\,\dd\mu_\theta.
\]
由定理 \ref{thm:xi-real-input-40-output-freezing-positive-entropy-halffill}，
\[
\lim_{\theta\to-\infty}P'(\theta)=0,
\qquad
\lim_{\theta\to+\infty}P'(\theta)=\frac12.
\]
令 \(\theta\to\pm\infty\) 即得
\[
\int \mathbf 1_{\mathrm{out}=1}\,\dd\mu_\infty^{-}=0,
\qquad
\int \mathbf 1_{\mathrm{out}=1}\,\dd\mu_\infty^{+}=\frac12.
\]
又同一定理已把 \(\mu_\infty^{+}\) 识别为由四态矩阵 \(B_\infty\) 的正熵地基层所支配的零温极限，故其饱和值严格停在 \(1/2\) 而非 \(1\)。证毕。
\end{proof}
\leanverified{paper\_xi\_real\_input\_40\_output\_freezing\_half\_filling\_empty\_phase}

\begin{theorem}[地基层熵、端点极点尺度与零温矩阵谱半径的平方律]\label{thm:xi-real-input-40-ground-entropy-endpoint-square-law}
沿用定理 \ref{thm:xi-real-input-40-output-freezing-positive-entropy-halffill} 的记号，并记
\[
z_\star(u):=\lambda(u)^{-1},
\qquad
b:=\rho(B_\infty)^{-1}.
\]
则有精确恒等式
\[
\boxed{\;
h_{\mathrm{ground}}
=
\log c
=
-\frac12\log b
=
\frac12\log\rho(B_\infty).\;}
\]
并且当 \(u\to+\infty\) 时，
\[
\boxed{\;
u\,z_\star(u)^2\to b=e^{-2h_{\mathrm{ground}}},
\qquad
z_\star(u)\sim e^{-h_{\mathrm{ground}}}u^{-1/2},
\qquad
\lambda(u)\sim e^{h_{\mathrm{ground}}}u^{1/2}.\;}
\]
\end{theorem}
\begin{proof}
由定理 \ref{thm:xi-real-input-40-output-freezing-positive-entropy-halffill}，
\[
h_{\mathrm{ground}}=\log c=\frac12\log\rho(B_\infty).
\]
故
\[
b=\rho(B_\infty)^{-1}=c^{-2}=e^{-2h_{\mathrm{ground}}},
\]
即得第一条盒装恒等式。

另一方面，命题 \ref{prop:real-input-40-pressure-freezing} 给出
\[
\lambda(u)=c\,u^{1/2}(1+o(1))
\qquad(u\to+\infty).
\]
取倒数便得
\[
z_\star(u)=\lambda(u)^{-1}=c^{-1}u^{-1/2}(1+o(1))
=
e^{-h_{\mathrm{ground}}}u^{-1/2}(1+o(1)).
\]
于是
\[
u\,z_\star(u)^2\to c^{-2}=b=e^{-2h_{\mathrm{ground}}},
\]
再将上式取倒数即得
\[
\lambda(u)\sim e^{h_{\mathrm{ground}}}u^{1/2}.
\]
证毕。
\end{proof}

\endinput


\endinput


\endinput

\paragraph{escort 二态极限、相对熵与 Shannon 常数的闭式回收}
在 escort 温度轴的末位 Gibbs 极限、对数多重度两原子分布以及 Rényi--Shannon 常数谱都已固定之后，还可把 bin-fold 的实参数 escort 族进一步压缩为下列四条直接后果。它们把末位边际、缩放退化度、对数退化度的一二阶矩、相对于稳定层均匀基线的 KL 缺口以及 Shannon 熵常数统一归结为单一 Bernoulli 参数
\[
\theta(a):=\frac{1}{1+\varphi^{a+1}}
\qquad (a\ge 0).
\]

\begin{theorem}[escort 纤维测度的二态极限律]\label{thm:xi-time-part54ab-escort-two-state-limit-law}
\leanverified{paper\_xi\_time\_part54ab\_escort\_two\_state\_limit\_law}
对任意实数 \(a\ge 0\)，定义 escort 纤维测度
\[
\pi_{m,a}(w):=\frac{d_m^{\mathrm{bin}}(w)^a}{S_a^{\mathrm{bin}}(m)},
\qquad
S_a^{\mathrm{bin}}(m):=\sum_{u\in X_m}d_m^{\mathrm{bin}}(u)^a.
\]
令 \(W_m^{(a)}\sim \pi_{m,a}\)，并记
\[
p_1(a):=\theta(a)=\frac{1}{1+\varphi^{a+1}},
\qquad
p_0(a):=1-\theta(a)=\frac{\varphi^{a+1}}{1+\varphi^{a+1}}.
\]
则有
\[
\mathbb P\bigl(W_m^{(a)}(m)=1\bigr)\longrightarrow p_1(a),
\qquad
\mathbb P\bigl(W_m^{(a)}(m)=0\bigr)\longrightarrow p_0(a).
\]
进一步，若定义缩放退化度
\[
Y_m^{(a)}:=\Bigl(\frac{\varphi}{2}\Bigr)^m d_m^{\mathrm{bin}}(W_m^{(a)}),
\]
则
\[
Y_m^{(a)}\Longrightarrow Y_\infty^{(a)},
\qquad
Y_\infty^{(a)}=
\begin{cases}
1,& \text{概率 }p_0(a),\\[4pt]
\varphi^{-1},& \text{概率 }p_1(a).
\end{cases}
\]
\end{theorem}
\begin{proof}
末位边际极限正是定理 \ref{thm:xi-fold-escort-lastbit-gibbs-limit} 在参数 \(q=a\) 处的结论。

再记
\[
X_m^{(a)}:=\log d_m^{\mathrm{bin}}(W_m^{(a)})-m\log\!\Bigl(\frac{2}{\varphi}\Bigr).
\]
则
\[
X_m^{(a)}=\log Y_m^{(a)}.
\]
由定理 \ref{thm:xi-fold-escort-log-multiplicity-two-atom} 在 \(q=a\) 处的结论，
\[
X_m^{(a)}\Longrightarrow
p_0(a)\delta_0+p_1(a)\delta_{-\log\varphi}.
\]
对上式应用连续映射 \(x\mapsto e^x\)，便得到
\[
Y_m^{(a)}\Longrightarrow
p_0(a)\delta_1+p_1(a)\delta_{\varphi^{-1}},
\]
即所述二点极限律。证毕。
\end{proof}

\leanverified{paper\_xi\_time\_part54ab\_escort\_logdefect\_moments}
\begin{theorem}[escort 对数退化度的一阶矩与方差闭式]\label{thm:xi-time-part54ab-escort-logdefect-moments}
沿用定理 \ref{thm:xi-time-part54ab-escort-two-state-limit-law} 的记号。则当 \(m\to\infty\) 时有
\[
\mathbb E_{\pi_{m,a}}\!\bigl[\log d_m^{\mathrm{bin}}(W_m^{(a)})\bigr]
=
m\log\!\Bigl(\frac{2}{\varphi}\Bigr)
-\frac{\log\varphi}{1+\varphi^{a+1}}
+O\!\Bigl(\Bigl(\frac{\varphi}{2}\Bigr)^m\Bigr),
\]
以及
\[
\operatorname{Var}_{\pi_{m,a}}\!\bigl(\log d_m^{\mathrm{bin}}(W_m^{(a)})\bigr)
=
\frac{\varphi^{a+1}}{(1+\varphi^{a+1})^2}(\log\varphi)^2
+O\!\Bigl(\Bigl(\frac{\varphi}{2}\Bigr)^m\Bigr).
\]
\end{theorem}
\begin{proof}
第一式正是命题 \ref{prop:fold-bin-escort-lastbit} 在 \(q=a\) 处的一阶矩展开。

再记
\[
\varepsilon_m:=\Bigl(\frac{\varphi}{2}\Bigr)^m.
\]
由定理 \ref{thm:fold-bin-two-state-asymptotic} 的一致对数展开，
\[
\log d_m^{\mathrm{bin}}(w)
=
m\log\!\Bigl(\frac{2}{\varphi}\Bigr)-w_m\log\varphi+\rho_m(w),
\qquad
\sup_{w\in X_m}|\rho_m(w)|=O(\varepsilon_m).
\]
因此
\[
\log d_m^{\mathrm{bin}}(W_m^{(a)})
=
m\log\!\Bigl(\frac{2}{\varphi}\Bigr)
-(\log\varphi)\,W_m^{(a)}(m)
+O_{L^\infty}(\varepsilon_m).
\]
常数平移不影响方差，故
\[
\operatorname{Var}_{\pi_{m,a}}\!\bigl(\log d_m^{\mathrm{bin}}(W_m^{(a)})\bigr)
=
(\log\varphi)^2\operatorname{Var}_{\pi_{m,a}}\!\bigl(W_m^{(a)}(m)\bigr)
+O(\varepsilon_m).
\]
另一方面，命题 \ref{prop:fold-bin-escort-lastbit} 给出
\[
\mathbb P\bigl(W_m^{(a)}(m)=1\bigr)
=
\theta(a)+O(\varepsilon_m).
\]
由于 \(W_m^{(a)}(m)\) 是 Bernoulli 变量，便有
\[
\operatorname{Var}_{\pi_{m,a}}\!\bigl(W_m^{(a)}(m)\bigr)
=
\theta(a)\bigl(1-\theta(a)\bigr)+O(\varepsilon_m).
\]
代入并用
\[
\theta(a)\bigl(1-\theta(a)\bigr)
=
\frac{\varphi^{a+1}}{(1+\varphi^{a+1})^2},
\]
即得所示方差公式。证毕。
\end{proof}

\begin{theorem}[escort 相对于稳定层均匀基线的 KL 极限与 Bernoulli 约化]\label{thm:xi-time-part54ab-escort-uniform-kl-bernoulli-reduction}
\leanverified{paper\_xi\_time\_part54ab\_escort\_uniform\_kl\_bernoulli\_reduction}
设 \(u_m\) 为 \(X_m\) 上的均匀分布，
\[
u_m(w)=\frac{1}{|X_m|}=\frac{1}{F_{m+2}}.
\]
则对任意 \(a\ge 0\)，有
\[
D_{\mathrm{KL}}(\pi_{m,a}\Vert u_m)
=
D_{\mathrm{KL}}\!\Bigl(\mathrm{Bernoulli}(p_1(a))\Vert \mathrm{Bernoulli}(\varphi^{-2})\Bigr)
+O\!\Bigl(\Bigl(\frac{\varphi}{2}\Bigr)^m\Bigr).
\]
等价地，
\[
\lim_{m\to\infty}D_{\mathrm{KL}}(\pi_{m,a}\Vert u_m)
=
p_0(a)\log\!\bigl(p_0(a)\varphi\bigr)
+
p_1(a)\log\!\bigl(p_1(a)\varphi^2\bigr).
\]
\end{theorem}
\begin{proof}
由推论 \ref{cor:xi-fold-escort-uniform-kl-temperature-law} 在 \(q=a\) 处的公式，
\[
D_{\mathrm{KL}}(\pi_{m,a}\Vert u_m)
=
D(a)+O\!\Bigl(\Bigl(\frac{\varphi}{2}\Bigr)^m\Bigr),
\]
其中
\[
D(a)=
2\log\varphi-\log(\varphi+\varphi^{-a})
-\frac{a\log\varphi}{1+\varphi^{a+1}}.
\]
再用
\[
p_1(a)=\frac{1}{1+\varphi^{a+1}},
\qquad
p_0(a)=\frac{\varphi^{a+1}}{1+\varphi^{a+1}},
\qquad
1-\varphi^{-2}=\varphi^{-1},
\]
直接整理可得
\[
D(a)
=
p_0(a)\log\!\frac{p_0(a)}{\varphi^{-1}}
+
p_1(a)\log\!\frac{p_1(a)}{\varphi^{-2}},
\]
亦即
\[
D(a)=
D_{\mathrm{KL}}\!\Bigl(\mathrm{Bernoulli}(p_1(a))\Vert \mathrm{Bernoulli}(\varphi^{-2})\Bigr).
\]
这就得到第一式。第二式只是将同一 Bernoulli KL 按显式坐标写开。证毕。
\end{proof}

\begin{theorem}[escort 测度在稳定层上的 Shannon 熵闭式]\label{thm:xi-time-part54ab-escort-shannon-closed-form}
\leanverified{paper\_xi\_time\_part54ab\_escort\_shannon\_closed\_form}
沿用定理 \ref{thm:xi-time-part54ab-escort-uniform-kl-bernoulli-reduction} 的记号。则 escort 测度 \(\pi_{m,a}\) 的 Shannon 熵满足
\[
H(\pi_{m,a})
=
m\log\varphi-\frac12\log 5
+\frac{\varphi^{a+1}}{1+\varphi^{a+1}}\log\varphi
+h\!\left(\frac{1}{1+\varphi^{a+1}}\right)
+O\!\Bigl(\Bigl(\frac{\varphi}{2}\Bigr)^m\Bigr),
\]
其中
\[
h(p):=-p\log p-(1-p)\log(1-p).
\]
等价地，
\[
H(\pi_{m,a})
=
\log F_{m+2}
-D_{\mathrm{KL}}\!\Bigl(\mathrm{Bernoulli}(p_1(a))\Vert \mathrm{Bernoulli}(\varphi^{-2})\Bigr)
+O\!\Bigl(\Bigl(\frac{\varphi}{2}\Bigr)^m\Bigr).
\]
\end{theorem}
\begin{proof}
由均匀基线的恒等式，
\[
H(\pi_{m,a})=\log|X_m|-D_{\mathrm{KL}}(\pi_{m,a}\Vert u_m)=\log F_{m+2}-D_{\mathrm{KL}}(\pi_{m,a}\Vert u_m).
\]
结合定理 \ref{thm:xi-time-part54ab-escort-uniform-kl-bernoulli-reduction}，立即得到第二个公式。

再由 Binet 渐近，
\[
\log F_{m+2}
=(m+2)\log\varphi-\frac12\log 5+O(\varphi^{-2m}).
\]
将
\[
D_{\mathrm{KL}}\!\Bigl(\mathrm{Bernoulli}(p_1(a))\Vert \mathrm{Bernoulli}(\varphi^{-2})\Bigr)
=
p_0(a)\log\!\bigl(p_0(a)\varphi\bigr)
+
p_1(a)\log\!\bigl(p_1(a)\varphi^2\bigr)
\]
代入，并用 \(p_0(a)+p_1(a)=1\) 整理，得到
\[
\begin{aligned}
H(\pi_{m,a})
&=
(m+2)\log\varphi-\frac12\log 5
-p_0(a)\log\!\bigl(p_0(a)\varphi\bigr)
-p_1(a)\log\!\bigl(p_1(a)\varphi^2\bigr)
+O\!\Bigl(\Bigl(\frac{\varphi}{2}\Bigr)^m\Bigr)\\
&=
m\log\varphi-\frac12\log 5
+p_0(a)\log\varphi
-p_0(a)\log p_0(a)-p_1(a)\log p_1(a)
+O\!\Bigl(\Bigl(\frac{\varphi}{2}\Bigr)^m\Bigr).
\end{aligned}
\]
最后注意
\[
-p_0(a)\log p_0(a)-p_1(a)\log p_1(a)=h\bigl(p_1(a)\bigr),
\]
便得第一式。证毕。
\end{proof}

\endinput

\paragraph{自然扩张字母表极小性、整数椭圆模核共谱与分层 Shannon 预算}
沿用定理 \ref{thm:killo-natural-extension-branch-register}、定理
\ref{thm:killo-smith-loss-spectrum}、推论
\ref{cor:killo-kernel-size-general-modulus}、定理
\ref{thm:killo-ellipse-diagonal-prime-register-equivalence} 与定理
\ref{thm:xi-prime-slice-nontrivial-layer-exact-minimality} 的记号。前述链条已经分别给出自然扩张的分支指标寄存器实现、Smith 核谱的可逆读出、模数核大小的全局闭式，以及整数轴对齐椭圆的 prime-register 语义。由此还可进一步抽出如下四条彼此咬合的后果：H.153 型寄存器实现的最小有限字母表定律、整数椭圆的 primewise 无序模核共谱分类、对应 Dirichlet 级数的 \(\zeta\)-主因子，以及分层外置的平均 Shannon 预算律。

\begin{theorem}[自然扩张的 H.153 型分支寄存器实现具有精确最小有限字母表]\label{thm:xi-natural-extension-h153-minimal-finite-alphabet}
\leanverified{paper\_xi\_natural\_extension\_h153\_minimal\_finite\_alphabet}
设 \(X\) 为可数集合，\(T:X\twoheadrightarrow X\) 为有限纤维满射，并记
\[
M(T):=\sup_{y\in X}\#T^{-1}(y)\in \NN\cup\{\infty\}.
\]
对任意整数 \(M\ge 1\)，下列两条等价：
\begin{enumerate}
  \item 存在大小为 \(M\) 的有限字母表 \(A\)，以及一族单射
  \[
  \iota_y:T^{-1}(y)\hookrightarrow A\qquad(y\in X),
  \]
  使得由
  \[
  \beta_A(x):=\iota_{T(x)}(x)
  \]
  诱导的 H.153 型编码
  \[
  \mathcal C_A:X^{\leftarrow}\longrightarrow X\times A^{\NN},
  \qquad
  \mathcal C_A((x_0,x_{-1},x_{-2},\dots))
  :=
  \bigl(x_0,(\beta_A(x_{-1}),\beta_A(x_{-2}),\dots)\bigr)
  \]
  为单射，并把自然扩张左移共轭到右移插入系统
  \[
  \widehat T_A\bigl(x,(a_1,a_2,\dots)\bigr)
  :=
  \bigl(T(x),(\beta_A(x),a_1,a_2,\dots)\bigr).
  \]
  \item \(M\ge M(T)\)。
\end{enumerate}
特别地，若 \(M(T)<\infty\)，则 H.153 型自然扩张分支寄存器实现所需的最小有限字母表大小恰为 \(M(T)\)；若 \(M(T)=\infty\)，则不存在任何有限字母表实现。
\end{theorem}
\begin{proof}
先证必要性。若存在第 \(1\) 条所述实现，则对每个 \(y\in X\) 都已有单射
\[
\iota_y:T^{-1}(y)\hookrightarrow A.
\]
故
\[
\#T^{-1}(y)\le |A|=M
\qquad(\forall\,y\in X),
\]
从而
\[
M(T)=\sup_{y\in X}\#T^{-1}(y)\le M.
\]

再证充分性。若 \(M\ge M(T)\)，则对每个 \(y\in X\) 都有
\[
\#T^{-1}(y)\le M=|A|.
\]
因此可以逐点选取单射
\[
\iota_y:T^{-1}(y)\hookrightarrow A.
\]
由此定义 \(\beta_A\) 与 \(\mathcal C_A\)。定理 \ref{thm:killo-natural-extension-branch-register} 的递推反演证明逐字适用：给定
\[
\bigl(x_0,(a_1,a_2,\dots)\bigr)\in \mathcal C_A(X^{\leftarrow}),
\]
可沿像上反演逐步恢复唯一前史
\[
x_{-1}=\iota_{x_0}^{-1}(a_1),\qquad
x_{-2}=\iota_{x_{-1}}^{-1}(a_2),\ \dots,
\]
从而 \(\mathcal C_A\) 为单射；并且
\[
\mathcal C_A\circ \sigma=\widehat T_A\circ \mathcal C_A.
\]
故第 \(1\) 条成立。证毕。
\end{proof}

\begin{theorem}[整数轴对齐椭圆的模核谱只看见 primewise 无序轴指数]\label{thm:xi-integer-ellipse-mod-kernel-spectrum-primewise-unordered}
\leanverified{paper\_xi\_integer\_ellipse\_mod\_kernel\_spectrum\_primewise\_unordered}
沿用定理 \ref{thm:killo-ellipse-diagonal-prime-register-equivalence} 的椭圆记号。对
\[
E_{a,b},E_{c,d}\in \mathsf E_{\NN},
\qquad
A_{a,b}:=\operatorname{diag}(a,b),
\qquad
A_{c,d}:=\operatorname{diag}(c,d),
\]
并对每个素数 \(p\) 与整数 \(k\ge 1\) 定义
\[
K_p^{a,b}(k):=\#\ker(A_{a,b}\bmod p^k).
\]
则成立等价
\[
K_p^{a,b}(k)=K_p^{c,d}(k)\quad(\forall\,p,\ \forall\,k\ge 1)
\iff
\{v_p(a),v_p(b)\}=\{v_p(c),v_p(d)\}\quad(\forall\,p).
\]
亦即，两个有序整数轴对齐椭圆具有完全相同的模核谱，当且仅当每个素数上的两根轴指数只发生逐素数的坐标交换。

进一步，定义
\[
\omega_{\neq}(a,b):=\#\{\,p:\ v_p(a)\neq v_p(b)\,\}.
\]
则与 \(E_{a,b}\) 共谱的有序椭圆恰有
\[
2^{\omega_{\neq}(a,b)}
\]
个。
\end{theorem}
\begin{proof}
对对角矩阵
\[
A_{a,b}=\operatorname{diag}(a,b)
\]
作 Smith 归约时，其 Smith 不变量即为
\[
d_1=\gcd(a,b),
\qquad
d_2=\operatorname{lcm}(a,b).
\]
故对每个素数 \(p\)，若记
\[
\alpha_p:=v_p(a),
\qquad
\beta_p:=v_p(b),
\]
则
\[
v_p(d_1)=\min(\alpha_p,\beta_p),
\qquad
v_p(d_2)=\max(\alpha_p,\beta_p).
\]
由定理 \ref{thm:killo-smith-loss-spectrum}，模核谱
\[
\bigl(K_p^{a,b}(k)\bigr)_{p,k}
\]
可逆恢复每个 \(p\) 上的 Smith 赋值对
\[
\bigl(v_p(d_1),v_p(d_2)\bigr)
=
\bigl(\min(\alpha_p,\beta_p),\max(\alpha_p,\beta_p)\bigr).
\]
因此它所恢复的恰是无序二元组
\[
\{\alpha_p,\beta_p\}=\{v_p(a),v_p(b)\}.
\]
这就证明了前一条等价关系。

现计算共谱有序椭圆的数目。对每个满足 \(\alpha_p\neq \beta_p\) 的素数 \(p\)，保持无序集合
\[
\{\alpha_p,\beta_p\}
\]
不变时，只有两种可能的有序分配：把较小指数放在第一轴、较大指数放在第二轴，或反之交换。若 \(\alpha_p=\beta_p\)，则交换不产生新对象。由于不同素数上的选择彼此独立，故总共有
\[
2^{\omega_{\neq}(a,b)}
\]
个有序椭圆与 \(E_{a,b}\) 共谱。证毕。
\end{proof}

\begin{corollary}[整数轴对齐椭圆的模核 Dirichlet 级数是 \texorpdfstring{$\zeta$}{zeta} 的有限 Euler 修正]\label{cor:xi-integer-ellipse-mod-kernel-dirichlet-zeta}\leanverified{paper\_xi\_integer\_ellipse\_mod\_kernel\_dirichlet\_zeta}
沿用定理 \ref{thm:xi-integer-ellipse-mod-kernel-spectrum-primewise-unordered} 的记号。对每个 \(N\ge 1\)，定义模数核计数函数
\[
K_{a,b}(N):=\#\ker(A_{a,b}\bmod N).
\]
则有精确闭式
\[
K_{a,b}(N)=\gcd(a,N)\gcd(b,N).
\]
因此 Dirichlet 级数
\[
\mathcal Z_{a,b}(s):=\sum_{N\ge 1}\frac{K_{a,b}(N)}{N^s}
\]
在 \(\Re(s)>1\) 上绝对收敛，并具有 Euler 乘积分解
\[
\mathcal Z_{a,b}(s)
=
\prod_p
\left(
\sum_{k\ge 0}p^{-ks+\min(k,\alpha_p)+\min(k,\beta_p)}
\right),
\qquad
\alpha_p:=v_p(a),\ \beta_p:=v_p(b).
\]
更进一步，
\[
\mathcal Z_{a,b}(s)=\zeta(s)\,R_{a,b}(s),
\]
其中
\[
R_{a,b}(s):=
\prod_{p\mid ab}
\Bigl(1-p^{-s}\Bigr)
\sum_{k\ge 0}p^{-ks+\min(k,\alpha_p)+\min(k,\beta_p)}
\]
是仅在坏素数 \(p\mid ab\) 上出现的有限 Euler 修正。因而 \(\mathcal Z_{a,b}(s)\) 在 \(s=1\) 处具有唯一简单极点。

此外，
\[
\mathcal Z_{a,b}(s)=\mathcal Z_{c,d}(s)
\iff
\bigl(K_p^{a,b}(k)\bigr)_{p,k}=\bigl(K_p^{c,d}(k)\bigr)_{p,k},
\]
从而也等价于 \(E_{a,b}\) 与 \(E_{c,d}\) 在定理 \ref{thm:xi-integer-ellipse-mod-kernel-spectrum-primewise-unordered} 的意义下共谱。
\end{corollary}
\begin{proof}
由推论 \ref{cor:killo-kernel-size-general-modulus}，对任意 \(A\in M_{m\times n}(\ZZ)\) 都有
\[
\#\ker(A_N)=N^{\,n-m}\prod_{i=1}^{m}\gcd(d_i,N),
\]
其中 \(d_i\) 为 Smith 不变量。对
\[
A_{a,b}=\operatorname{diag}(a,b)
\]
有 \(m=n=2\)，且 Smith 不变量为
\[
d_1=\gcd(a,b),
\qquad
d_2=\operatorname{lcm}(a,b).
\]
故
\[
K_{a,b}(N)=\gcd(d_1,N)\gcd(d_2,N).
\]
逐素数比较即可验证
\[
\gcd(d_1,N)\gcd(d_2,N)=\gcd(a,N)\gcd(b,N),
\]
于是得到所示精确闭式。

由右端关于 \(N\) 的积性，Dirichlet 级数分解为 Euler 乘积
\[
\mathcal Z_{a,b}(s)
=
\prod_p \sum_{k\ge 0}\frac{K_{a,b}(p^k)}{p^{ks}}.
\]
而
\[
K_{a,b}(p^k)=p^{\min(k,\alpha_p)+\min(k,\beta_p)},
\]
故得到所示局部因子公式。

若 \(p\nmid ab\)，则 \(\alpha_p=\beta_p=0\)，从而
\[
\sum_{k\ge 0}p^{-ks+\min(k,\alpha_p)+\min(k,\beta_p)}
=
\sum_{k\ge 0}p^{-ks}
=
\frac{1}{1-p^{-s}}.
\]
把全部好素数因子提出，便得到
\[
\mathcal Z_{a,b}(s)
=
\zeta(s)
\prod_{p\mid ab}
\Bigl(1-p^{-s}\Bigr)
\sum_{k\ge 0}p^{-ks+\min(k,\alpha_p)+\min(k,\beta_p)}.
\]
由于坏素数只有有限多个，\(R_{a,b}(s)\) 为有限 Euler 乘积；并且每个局部修正因子都是 \(p^{-s}\) 的有理函数，故 \(R_{a,b}\) 在 \(s=1\) 的邻域全纯。又因各局部因子在 \(s=1\) 处都取正实值，故 \(R_{a,b}(1)\neq 0\)。因此 \(\mathcal Z_{a,b}(s)\) 只在 \(s=1\) 继承 \(\zeta(s)\) 的唯一简单极点。

最后，若 \(\mathcal Z_{a,b}(s)=\mathcal Z_{c,d}(s)\)，则两边 Dirichlet 系数逐项相同，特别地对每个素数幂 \(N=p^k\) 有
\[
K_p^{a,b}(k)=K_p^{c,d}(k).
\]
反向蕴含显然。再结合定理 \ref{thm:xi-integer-ellipse-mod-kernel-spectrum-primewise-unordered}，便得到与共谱条件的等价性。证毕。
\end{proof}

\begin{theorem}[分层外置的平均 Shannon 极小预算由层间条件熵完全控制]\label{thm:xi-layered-externalization-average-shannon-budget}
\leanverified{paper\_xi\_layered\_externalization\_average\_shannon\_budget}
设
\[
X_0\xrightarrow{\pi_0}X_1\xrightarrow{\pi_1}\cdots\xrightarrow{\pi_{k-1}}X_k
\]
为有限纤维满射塔，并在 \(X_0\) 上赋予任意概率分布 \(\mu\)。记
\[
Y_0\sim\mu,
\qquad
Y_{j+1}:=\pi_j(Y_j)\qquad(0\le j\le k-1),
\]
并对每个 \(j\) 与 \(y\in X_{j+1}\) 记纤维
\[
F_j(y):=\pi_j^{-1}(y).
\]

对每个 \(j\) 与 \(y\in X_{j+1}\)，取一个二进制前缀码
\[
c_{j,y}:F_j(y)\longrightarrow \{0,1\}^{\ast}.
\]
把分层侧信息定义为
\[
\Psi_c(x_0)
:=
\Bigl(
x_k,\,
c_{k-1,x_k}(x_{k-1})\cdots c_{1,x_2}(x_1)c_{0,x_1}(x_0)
\Bigr),
\]
其中 \(x_{j+1}=\pi_j(x_j)\)。记平均侧信息长度为
\[
L(c):=
\EE_\mu\!\left[
\sum_{j=0}^{k-1}
\bigl|c_{j,Y_{j+1}}(Y_j)\bigr|
\right].
\]
则有下界
\[
L(c)\ge
\frac{1}{\log 2}\sum_{j=0}^{k-1}H_\mu(Y_j\mid Y_{j+1}).
\]
并且存在一族二进制前缀码 \(\{c_{j,y}\}\) 使
\[
L(c)\le
\frac{1}{\log 2}\sum_{j=0}^{k-1}H_\mu(Y_j\mid Y_{j+1})+k.
\]
若上述条件熵总和有限，则由链式法则右端也可写为
\[
\frac{1}{\log 2}H_\mu(Y_0\mid Y_k)
\quad\text{与}\quad
\frac{1}{\log 2}H_\mu(Y_0\mid Y_k)+k.
\]
\end{theorem}
\begin{proof}
先证下界。固定一层 \(j\) 与一点 \(y\in X_{j+1}\)。由于
\[
c_{j,y}:F_j(y)\to\{0,1\}^{\ast}
\]
是二进制前缀码，Shannon 下界给出其在条件分布
\[
\mu(\cdot\mid Y_{j+1}=y)
\]
下的平均码长满足
\[
\sum_{x\in F_j(y)}
\mu(x\mid y)\,|c_{j,y}(x)|
\ge
\frac{1}{\log 2}
\left(
-\sum_{x\in F_j(y)}\mu(x\mid y)\log\mu(x\mid y)
\right).
\]
右端恰为
\[
\frac{1}{\log 2}H_\mu(Y_j\mid Y_{j+1}=y).
\]
再对 \(y\) 以 \(Y_{j+1}\) 的分布加权，即得
\[
\EE_\mu\bigl[|c_{j,Y_{j+1}}(Y_j)|\bigr]
\ge
\frac{1}{\log 2}H_\mu(Y_j\mid Y_{j+1}).
\]
对 \(j=0,\dots,k-1\) 求和，便得到
\[
L(c)\ge
\frac{1}{\log 2}\sum_{j=0}^{k-1}H_\mu(Y_j\mid Y_{j+1}).
\]

再证上界。对每个固定的 \((j,y)\)，在有限集 \(F_j(y)\) 上按条件分布
\[
\mu(\cdot\mid Y_{j+1}=y)
\]
构造 Shannon--Fano 码；若某些点具有零条件概率，则在有限纤维上任意补全为前缀码即可，而不会影响条件平均长度。于是可取 \(c_{j,y}\) 使
\[
\sum_{x\in F_j(y)}
\mu(x\mid y)\,|c_{j,y}(x)|
\le
\frac{1}{\log 2}H_\mu(Y_j\mid Y_{j+1}=y)+1.
\]
再对 \(y\) 加权，得到
\[
\EE_\mu\bigl[|c_{j,Y_{j+1}}(Y_j)|\bigr]
\le
\frac{1}{\log 2}H_\mu(Y_j\mid Y_{j+1})+1.
\]
对全部层求和，便有
\[
L(c)\le
\frac{1}{\log 2}\sum_{j=0}^{k-1}H_\mu(Y_j\mid Y_{j+1})+k.
\]

最后，若上述熵量皆有限，则由于 \(Y_{j+1}\) 是 \(Y_j\) 的确定函数，
\[
H_\mu(Y_j,Y_{j+1})=H_\mu(Y_j),
\]
从而
\[
H_\mu(Y_j\mid Y_{j+1})=H_\mu(Y_j)-H_\mu(Y_{j+1}).
\]
望远镜求和即得
\[
\sum_{j=0}^{k-1}H_\mu(Y_j\mid Y_{j+1})
=
H_\mu(Y_0)-H_\mu(Y_k)
=
H_\mu(Y_0\mid Y_k).
\]
证毕。
\end{proof}

\noindent
这些结果表明：H.153 型自然扩张分支寄存器实现的有限字母表消耗并非可任意压缩，而是恰由最大纤维基数锁定；H.154--H.156 在整数轴对齐椭圆语义下并不恢复逐素数“哪根轴承载哪一指数”，而只恢复每个素数上的无序二元组；相应模核 Dirichlet 级数则正是 \(\zeta(s)\) 的有限 Euler 修正。另一方面，H.151--H.152 所给出的“每个活跃层至少一枚有效 prime-slice”最坏情形定律，在平均意义下精确退化为由层间条件熵总和控制的 Shannon 极小预算。这便把自然扩张、Smith 模核谱、整数椭圆与分层外置的硬/软复杂度进一步压缩为同一条可审计结构链。

\endinput


\endinput
